% Gerado por tools/tei2latex.py a partir de revisado.xml. Regenerar, não editar.

%%>>> parte-2

%%>>> cap-6

%%>>> art-38

\pg{443} CAPITULO SEGUNDO

\textbf{ESSENCIA} E \textbf{ATTRIBUTOS DE DEUS}

\textsc{SUMMARIO}: --- Essencia de \textsc{Deus}, seu constitutivo e attributos. --- Attributos

de \textsc{Deus} negativos. --- Attributos de \textsc{Deus} positivos.

%%>>> art-39

\section*{ARTIGO I.\\ Essencia de Deus, seu constitutivo e attributos.}
\addcontentsline{toc}{section}{ARTIGO I. --- Essencia de Deus, seu constitutivo e attributos.}

22. \textbf{Essencia de Deus}. --- \emph{Essencia} de \textsc{Deus} é o que faz com que Elle seja \textsc{Deus}, e se distinga de todos os outros seres.

23. \textbf{Essencia physica} e metaphysica de \textbf{Deus}. --- A essencia \emph{physica} de \textsc{Deus} é constituida por todas as perfeições, de que a Divindade é dotada. A \emph{metaphysica} é a perfeição fundamental, que \emph{se concebe} como raiz e fonte de todas as outras perfeições divinas, e pela qual \textsc{Deus} se distingue das creaturas. --- A essencia \emph{physica} representa \textsc{Deus} como é em si mesmo; a \emph{metaphysica} exprime \textsc{Deus} como é em relação ao nosso modo de entender \footnote{Lembramos que a essencia --- é \emph{physica}, quando os seus elementos constitutivos são \emph{reaes} e \emph{concretos} (assim o \emph{corpo organico} e a \emph{alma racional} constituem a essencia \emph{physica} do homem), --- é \emph{metaphysica}, quando os seus elementos constitutivos são \emph{logicos} e \emph{abstractos}, denotando o \emph{genero} e a \emph{differença} (assim a \emph{animalidade} e a \emph{racionalidade} constituem a essencia \emph{metaphysica} do homem). --- Da essencia brotam as \emph{propriedades}. Estas distinguem-se \emph{realmente} da essencia \emph{physica}, porque não a constituem mas resultam d'ella (assim a intelligencia e a vontade são realmente distinctas da essencia \emph{physica} da nossa alma), --- mas podem não distinguir-se \emph{realmente} da essencia \emph{metaphysica}, porque, ainda que não possam considerar-se como elementos \emph{constitutivos} da essencia, comtudo podem identificar-se com esses elementos (assim a \emph{immortalidade}, embora não possa considerar-se como um elemento constitutivo da essencia \emph{metaphysica} da nossa alma, comtudo não se distingue \emph{realmente} da mesma essencia). E' claro que não podemos fallar da essencia de \textsc{Deus}, como fallamos da essencia das creaturas. Em Deus, a essencia \emph{physica} significa e abrange toda a perfeição, que compete ao Ente infinito, sem sombra de imperfeição, --- e a essencia \emph{metaphysica} exclue toda e qualquer composição. Por isso o \emph{constitutivo} da essencia \emph{metaphysica} de Deus não é uma coisa realmente distincta das outras propriedades, e a sua relação com estas deve, mais ou menos, conceber-se como se concebe a relação do ente, considerado na sua razão communissima, com as suas propriedades \emph{transcendentaes}.\par Dizemos que a essencia \emph{metaphysica} exprime \textsc{Deus} como Elle é em \emph{relação} ao \emph{nosso modo de entender}. Por quanto, considerando os seres creados,\pgn{444}encontramos em cada um d'elles uma perfeição, que é a primeira a ser conhecida, que constitue o mesmo ser numa determinada especie, distinguindo-o de todos os outros, e que é a raiz e a fonte de todas as outras perfeições. E, como subimos ao conhecimento de \textsc{Deus} por meio das creaturas, procuramos encontrar em \textsc{Deus} essa perfeição \emph{primeira}, \emph{caracteristica e radical}; embora na essencia de \textsc{Deus}, como é em si mesma, não se encontre attributo algum, que seja anterior aos outros attributos e de que estes fluam. --- Este modo de entender a essencia de \textsc{Deus} é o resultado da limitação da nossa intelligencia; não podemos entender o infinito senão por partes, e, ainda assim, com muita imperfeição. Todavia, a propria intelligencia emenda ou explica o seu modo de ver; pois que, depois de ter considerado um attributo divino, como fundamento dos outros, prova que, em \textsc{Deus}, não só não ha distincção alguma entre a essencia e os attributos, mas nem sequer entre os mesmos attributos.}.

\pg{444} 24. \textbf{A essencia physica de Deus consiste na infinidade.}

a) Em Deus, o ser (ou a existencia) não é recebido nem subsiste na essencia, mas é \emph{subsistente por si}; visto que o ser de Deus não se distingue da sua essencia, e o que não se distingue de uma coisa, não pode ser recebido nem subsistir nella. Ora o ente, cujo ser não é recebido nem subsiste na essencia, é necessariamente infinito; porque o ser é limitado ou finito só porque e quando é recebido numa certa e determinada essencia. Logo Deus é infinito, isto é, a essencia \emph{physica} de \textsc{Deus} consiste na \emph{infinidade} \footnote{Desenvolvamos este argumento. Uma substancia, como dissemos na \emph{Ontologia}, pode \emph{subsistir}, ou \emph{existir em si}, por dois modos: ou porque é a existencia, ou porque \emph{tem} a existencia. No primeiro caso, a existencia (ou o ser) não é recebida na essencia, ou natureza, como na sua potencia proporcionada, e por isso é acto puro e subsistente por si; --- no segundo, a existencia (ou o ser) é recebida na essencia, como acto e principio de subsistencia, e por isso não é acto puro, nem subsiste por si, mas subsiste na essencia, como na sua potencia. Ora, em Deus, a essencia e o ser, ou a existencia, não se distinguem \emph{realmente}, porque o ser não é accidental, mas é essencial á essencia divina; porque esta é \emph{essencialmente} existente. Se o ser não se distingue \emph{realmente} da essencia divina, Deus, assim como é a sua essencia, isto é, a sua Divindade, assim tambem é o seu proprio ser, ou a sua propria existencia. Mas o ser, que não se distingue da essencia, não é recebido na propria essencia; porque o ser, que é recebido na essencia, deve necessariamente distinguir-se d'esta, como o acto se distingue necessariamente da potencia (assim a côr, recebida numa parede, distingue-se da propria parede). O ser, que não é recebido numa essencia, não é limitado ou determinado a uma certa especie; visto que o ser emtanto se limita ou se determina a uma certa especie, emquanto é recebido numa certa potencia ou natureza (assim o ser \emph{humano} é limitado ou determinado a uma certa especie, e não a uma outra, porque é recebido na \emph{essencia} ou \emph{natureza humana}). Portanto, o ser, que não é recebido numa essencia, não é limitado ou circumscripto, mas é illimitado ou infinito. E, como o \emph{ser}, por ser \emph{acto},\pgn{445}indica e é \emph{perfeição}, segue-se que o ser subsistente por si é a perfeição subsistente, illimitada, infinita; de modo que não ha nem pode conceber-se entitade, perfeição, attributo, que esse ser não possua, assim como a alvura, se fosse subsistente em si, havia de conter tudo o que tem razão de alvor. Deus, portanto, sendo o ser subsistente por si, é infinito na sua perfeição, E' o raciocinio do Doutor Angelico, quando diz: « Deus est ipsum esse per se subsistens; ex quo oportet quod totam perfectionem essendi in se contineat. Manifestum est enim quod si aliquod calidum non habeat totam perfectionem calidi, hoc ideo est, quia calor non participatur secundum perfectam rationem; sed, si calor esset per se subsistens, non posset ei aliquid deesse de virtute caloris. Unde cum Deus sit ipsum esse subsistens, nihil de perfectione essendi potest ei deesse. Omnium autem perfectiones pertinent ad perfectionem essendi; secundum hoc enim aliqua perfecta sunt quod aliquo modo esse habeant. Unde sequitur quod nullius rei perfectio Deo desit » (\emph{Sum. Theolog.}, p. I, q. 4, a. 2). --- Mais adiante accrescenta: « Cum esse divinum non sit esse receptum in aliquo; sed ipse sit suum esse subsistens, manifestum est quod ipse Deus sit infinitus et perfectus » (\emph{Ib.}, q. 7, a. 1). E noutro logar: « Deus est infinitus: quod sic videri potest. Esse enim hominis terminatum est ad hominis speciem, quia est receptum in natura speciei humanae, et simile est de esse equi, vel cujuslibet creaturae. Esse autem Dei, quum non sit in aliquo receptum, sed sit esse purum, non limitatur ad aliquem modum perfectionis essendi, sed totum esse in se habet » (\emph{De Pot.}, q. 1, a. 2).}.

\pg{445} b) A essencia \emph{physica} de \textsc{Deus} é constituida por todas as perfeições, de que a Divindade é dotada. Ora as perfeicoes, de que a Divindade é dotada, são \emph{infinitas} ; visto que \textsc{Deus}, sendo \emph{immutavel} e \emph{necessario}, deve possuir todas as possiveis perfeições; pois, se Lhe faltasse alguma, seria, com relação a esta, \emph{mutavel} e \emph{contingente}, \emph{e} por isso dependeria de uma causa. Logo á essencia \emph{physica} de \textsc{Deus} consiste na \emph{infinidade}.

c) Se as perfeicoes, que constituem a essencia \emph{physica} de \textsc{Deus}, fossem \emph{finitas} ou \emph{limitadas}, essa limitação deveria provir --- ou de uma causa \emph{extrinseca}, --- ou da indole da propria essencia divina, em que as perfeicoes são recebidas. Ora não pode provir --- nem de uma causa \emph{eztrinseca}, porque Dus, sendo a Causa \emph{primeira}, não depende de outra causa, --- nem da indole da propria essencia divina, porque esta, sendo perfeita por \emph{si mesma}, e não por \emph{participação}, não pode deixar de possuir \emph{a} plenitude da perfeição. Logo as perfeições, que constituem a essencia \emph{physica} de \textsc{Deus}, são \emph{infinitas}. Logo a essencia \emph{physica} de \textsc{Deus} consiste na \emph{infinidade} \footnote{A \emph{infinidade} de \textsc{Deus} \emph{indirectamente} demonstra-se pelo facto de se não encontrar nº Elle nenhuma das causas, que limitam as creaturas. Na verdade, as creaturas, que são compostas de \emph{materia} e \emph{forma}, são limitadas, porque á \emph{forma} limita á \emph{materia}, quando a determina a constituir uma certa especie e não outra, --- e a \emph{materia} limita a \emph{forma}, tornando-a \emph{sin}-,\pgn{446}\emph{gular} e \emph{individua}. Ora \textsc{Deus} nao pode ser limitado d'este modo: porque, como provaremos d'aqui a pouco, não é composto da \emph{materia} e \emph{forma}. --- As creaturas todas pertencem a um certo \emph{genero} e a uma certa \emph{especie}, e por isso são limitadas, porque não possuem as perfeições proprias dos ouires generos e das outras, especies. Ora \textsc{Deus}, sendo o ser \emph{necessario}, \emph{immutavel} e \emph{independente}, transcende todos os generos e todas as especies. --- As creaturas todas são compostas de \emph{essencia} e de \emph{ser}, ou \emph{existencia}, e por isso são limitadas, porque o \emph{ser}, sendo recebido numa essencia \emph{finita}, como é a essencia da creatura, não pode deixar de ser \emph{finito}. Ora, como dissemos, o ser divino, \emph{existindo} por assim o exigir a sua \emph{essencia}, subsiste por \emph{si} e não é recebido numa essencia, que o limite. --- As creaturas são limitadas pelo \emph{espaço} e pelo \emph{tempo}. \textsc{DEus}, por ser \emph{immenso}, não é limitado pelo \emph{espaço}, e, por ser \emph{eterno}, não é determinado pelo \emph{tempo}. Logo \textsc{Deus} é \emph{infinito}. --- S. Thomaz prova a mesma verdade, dizendo que Deus é \emph{infinito}; porque a razão da limitação num ente deriva ou da \emph{causa} do proprio ente ou da \emph{essencia}, em que o ser é recebido, mas Deus não tem causa e o seu ser não é recebido na essencia, mas, como provámos, é subsistente por si. Eis as suas palavras: « Ipsum esse, absolute consideratum, infinitum est... Si igitur alieujus esse sit finitum, oportet quod limitetur esse illud per aliquid aliud, quod sit aliqualiter causa illius esse, vel \emph{receptivum} ejus. Sed esse divini non potest esse aliqua causa, quia ipse necesse est esse per seipsum: nec esse ejus est receptivum, quum ipse sit suum esse, Ergo esse suum est infinitum, et ipse infinitus » (C. \emph{Gent}., lib. I, c. 48).\par Os erros ácerca da \emph{infinidade} de \textsc{Deus} podem reduzir-se a tres.\par a) Uns escriptores, não entendendo por \emph{infinito} senão a \emph{quantidade}, disseram que \textsc{Deus} não é infinito, porque uma \emph{quantidade infinita} repugna.\par Esta opinião refuta-se, negando o \emph{supposto}. A \emph{infinidade}, que se attribue a \textsc{Deus}, não é de \emph{quantidade}, mas sim de \emph{perfeição}.\par \emph{b}) Outros, como \emph{Locke}, \emph{Hobbes}, \emph{Robinet}, dizem que Deus é \emph{infinito}, não \emph{absoluta} mas \emph{relativamente}; pois faltam-lhe muitas perfeições.\par Esta opinião é absurda; porque \textsc{DEus}, como provámos, é \emph{essencialmente} infinito nas suas adoraveis perfeições. j\par c) Outros, como \emph{Hegel}, seguido por \emph{Vacherot}, \emph{Renan} e \emph{Taine}, não fazendo \textbar{} distincção entre a idea meramente \emph{subjectiva} do ente em geral e o Ente necessario, nem entre o conceito de \emph{infinidade} e o Infinito real, dizem que \textsc{DEus} não é um ente \emph{real} nem \emph{pessoal}, mas é a \emph{categoria do ideal}.\par Como é evidente, estas affirmações encerram o \emph{pantheismo}. Um \textsc{Deus} ideal, formado por nós, não é \textsc{Deus}. --- Não nos demoraremos na refutação do \emph{systema} de \emph{Hegel}; pois o refutámos na \emph{Cosmologia geral}, tratando do \emph{pantheismo}. Aqui apenas accrescentaremos que as theorias de \emph{Hegel}, servilmente copiadas ou imitadas por muitos escriptores, são tidas, na propria Allemanha, naquella conta que merecem. \emph{Schopenhauer} diz --- que \emph{Hegel} u é uma \emph{cabeça mediocre},\pgn{447}\emph{que por todos os meios quiz passar por um grande philosopho} », --- e que a philosophia de \emph{Hegel}, como a de \emph{Fichte} e de \emph{Schelling}, \emph{é} «\emph{um} minimum \emph{de pensamento diluido in quinhentas paginas de phraseologia nauseabunda} n...\par Davus, sendo infinito nas suas perfeições, deve possuir todas as perfeições, espalhadas nas creaturas; pois a causa \emph{efficiente} deve possuir toda a perfeição, que se encontra no effeito. --- E' o processo da \emph{causalidade}.\par Mas esta posse é diversa, segundo a natureza das perfeições. Por quanto as perfeições podem ser \emph{absolutas} ou \emph{miwtas}. +\par \emph{Absolutas} são as perfeições, que no seu conceito não encerram imperfeição alguma, nem se oppõem a perfeições mais elevadas, como são a \emph{vida}, a \emph{belleza}, a \emph{bondade}, a \emph{sabedoria}, etc. Estas perfeições attribuem-se a \textsc{Deus} no proprio conceito, que exprimem, mas por um modo mais \emph{eminente}, porque em \textsc{Deus} se encontram por um modo infinitamente mais perfeito e sublime do modo, por que se encontram nas creaturas, pois \textsc{DEus} não é \emph{bom}, \emph{santo}, \emph{amavel}, como é o homem. Por isso os nomes, com que exprimimos as perfeições \emph{absolutas} de \textsc{Deus}, não são \emph{univocos} nem \emph{equivocos}, mas \emph{analogos}. (\emph{Sum}. \emph{Th}., p. I, q. 19, a. 5). --- E' o processo da \emph{eminencia}.\par \emph{Mixtas} são as perfeições, que, embora no seu conceito exprimam uma entidade positiva, todavia encerram alguma imperfeição e se oppõem a uma perfeição mais elevada; tal é a \emph{corporeidade}, que, embora signifique uma coisa positiva, todavia encerra alguma imperfeição (pois o \emph{corpo}, por ser \emph{composto}, \emph{é dependente}, \emph{potencial}, \emph{finito}, ete.) e se oppõe a uma perfeição mais elevada, que é a \emph{espiritualidade} (pois o \emph{espirito} é mais nobre que o \emph{corpo}). Estas perfeicgóes não se attribuem a \textsc{Deus} no seu proprio conceito, mas só \emph{eminentemente}, emquanto attribuimos a \textsc{Deus} toda a entidade, que essas perfeições encerram, e afastamos d'Elle toda a imperfeição, de que as perfeições estão cercadas. --- E' o processo da \emph{exclusão}.\par D'ahi se vê que o processo da \emph{exclusão} é, para nós, o mais natural, para nos elevarmos ao conhecimento de \textsc{Deus}. Por quanto, devemos afastar de \textsc{Deus} --- não só as perfeições creadas, que no seu proprio conceito indicam alguma imperfeição, --- mas tambem todo o limite de que estão cercadas as perfeições creadas, que no seu conceito não encerram alguma imperfeição. D'este modo mais sabemos o que \textsc{Deus} não é, do que o que Elle é.\par Todavia, este processo, --- além de manifestar.a perfeição infinita de \textsc{DEus}, pois o Ente, de que se afasta toda a imperfeição, é perfeitissimo, --- subminisbra-nos um verdadeiro conhecimento de Dus, pois tanto melhor conheceremos um ente, quanto mais o distinguirmos de todos os outros seres. --- \textsc{DEUS}, pois, não é absolutamente \emph{i5cognoscivel}, como pretende \emph{H}. \emph{Spencer}.}.

25. \emph{A} \textbf{essencia} metaphysica de \textbf{Deus} consiste na subsistencia do proprio ser. --- A essencia \emph{metaphysica} de \textsc{Deus} consiste naquella perfeição, --- que se concebe como raiz e fonte de todas as outras perfeicoes divinas, e pela qual \textsc{Deus} se distingue das creaturas. Ora a \emph{subsistencia do proprio ser} é a perfeição \pg{446} \pg{447}  --- que se concebe como raiz e fonte de todas as outras perfeições divinas, porque, como provámos, Deus possue todas as perfeições, possiveis e imaginaveis, pelo facto de ser o \emph{proprio ser subsistente}, --- e pela qual se distingue das creaturas, visto que Deus é o Ente infinito, porque \emph{é} o proprio ser subsistente, e as creaturas são finitas, porque \emph{têm} o ser recebido numa determinada essencia. Logo a essencia \emph{metaphysica} de \textsc{Deus} consiste na \emph{subsistencia do proprio ser} \footnote{A conclusão é certa, e é uma consequencia dos principios expostos no numero anterior. O ser divino, pelo facto de não ser recebido em nenhuma\pgn{448}essencia ou potencia, não tem limites, encerra uma perfeição absoluta sem sombra de imperfeição. Repetimos com S. Thomaz: « \emph{Ex hoc quod Deus est ipsum esse subsistens}, \emph{oportet quod totam perfectionem essendi in se contineat} » (\emph{Sum. Th.}, p. I, q. 4, a. 2). Sendo subsistente, o ser divino distingue-se das creaturas, que não \emph{são} o ser, mas \emph{têm} o ser, recebido numa determinada potencia ou essencia e por isso limitado a uma certa especie. Como o ser \emph{é acto}, o ser que não é recebido numa potencia, mas que é acto \emph{puro}, é necessariamente infinito ou illimitado; ao passo que o ser, que é recebido numa potencia, é limitado por ella. Por isso, não só pode dizer-se que o \emph{constitutivo} da essencia divina consiste em que \textsc{Deus} \emph{é o ser subsistente por si}, mas tambem pode dizer-se que esse \emph{constitutivo} consiste em que \textsc{Deus} é \emph{acto purissimo}, \emph{não recebido em nenhuma potencia}. As expressões são synonimas. --- Concluimos --- que a razão formal e intrinseca, pela qual \textsc{Deus} é unico, eterno, infinito, é porque Elle é o proprio ser subsistente, é acto purissimo, cuja essencia é constituida pelo proprio ser, --- e que, pelo contrario, a razão, pela qual as creaturas são multiplicaveis, imperfeitas, finitas, é porque têm o ser recebido na essencia, e por isso são compostas de potencia e de acto, de essencia e de ser.\par Sendo Deus o proprio ser subsistente, o nome mais proprio de Deus é \textsc{JEHOVAH}, que significa: Qui EST --- \textsc{AQUELLE QUE} É. Por quanto, os nomes significam a forma ou essencia das coisas. Ora, em Deus, a essencia identifica-se com o ser, isto é, o ser é a sua propria essencia, de modo que á pergunta: \emph{o que é Deus}? podemos responder: \emph{Deus é o proprio ser} --- \emph{Deus est ipsum esse}. Logo o nome \emph{Qui est} é muito proprio de Deus. Com effeito, de nenhuma creatura podemos dizer que ella é o proprio ser; porque o \emph{ser}, nas \emph{creaturas}, não é a \emph{essencia}, nem é um \emph{predicado essencial}, mas é distincto da essencia, embora seja recebido nella; e por isso os nomes, que significam as coisas creadas, exprimem a essencia, e não e ser, assim \emph{homem} exprime a essencia de um individuo, isto é, a \emph{humanidade}, mas não exprime o \emph{ser} ou a \emph{existencia} d'esse individuo, que poderia não existir. Escutemos S. Thomaz: « Cum in omni quod est, sit considerare \emph{quidditatem suam}, per quam subsistit in natura determinata, et \emph{esse suum}, per quod dicitur de eo quod est in actu, hoc nomen \emph{res} imponitur rei a quidditate sua, hoc nomen \emph{qui est} vel \emph{ens} imponitur ab ipso actu essendi. Cum autem ita sit quod in qualibet re creata, essentia sua differat a suo esse, res illa proprie denominatur a quidditate sua et non ab actu essendi, sicut homo ab humanitate. In Deo autem ipsum esse suum est sua quidditas; et ideo nomen quod sumitur ab esse, proprie nominat ipsum, et est proprium nomen ejus; sicut proprium nomen hominis quod sumitur a quidditate sua » (I \emph{Sent}, Dist. 8, q. 1, a. 1).\par E não pode deixar de ser assim. Quando digo: \emph{Pedro é homem}, attribuo a esse individuo duas coisas: o \emph{ser}, ou a \emph{existencia}, e a \emph{natureza humana}, isto é, attribuo-lhe o ser limitado e determinado a uma certa especie; mas quando digo: \emph{Deus é aquelle que é,} quero affirmar que Deus não é este ou aquelle ente, mas que é o ser, todo o ser, o proprio ser, o ser infinito, de modo que tudo o que tem razão de ser deve encontrar-se e se encontra\pgn{449}em Deus. Eis o sentido do nome \emph{qui est}, que o proprio Deus manifestou a Moyses. Lê-se no \emph{Exodo}; « Ait Moyses ad Deum: Ecce ego vadam ad filios Israel, et dicam eis: Deus patrum vestrorum misit me ad vos. Si dixérint mihi: quod est nomen ejus? quid dicam eis? Dixit Deus ad Moysen: \textsc{Ego sum qui sum}. Ait: Sic dices filiis Israel: Qui EST misit me ad vos » (Cap. III, vv. 13, 14). --- E' este um nome mais universal e menos determinado que os outros; mas por isso mesmo é o mais proprio. Sendo a essencia divina infinitamente superior á intelligencia creada, o nome mais proprio de Deus é o que melhor traduz essa transcendencia e indefinibilidade, e exclue as imperfeitas determinações das creaturas. Diz S. Thomaz: « Quia divina comprehendi a nobis non possunt, convenientius a nobis significantur per nomina communia, quae indefinite aliquid significant, quam per nomina specialia, quae definite rei species exprimunt; unde hoc nomen \emph{qui est}, quod secundum Damascenum significat substantiae pelagus infinitum, convenientissimum nomen dicitur esse » (\emph{De pot}., q. 10, a. 1 ad 9).\par Nem se diga que o nome \emph{qui est} não é o mais proprio de Deus, porque não o distingue das creaturas, pois Deus é e as creaturas tambem \emph{são}. Por quanto, ha uma grande differença entre o \emph{ser} de Deus e o \emph{ser} das creaturas. O ser de Deus, por ser subsistente em si, identifica-se com a essencia, e por isso se torna individuo e distincto de todos os outros entes; ao passo que o ser das creaturas, por não ser subsistente em si, não se identifica com a essencia, mas é recebido e subsiste na essencia, e por isso se torna individuo e distincto por causa da mesma essencia. Portanto o nome \emph{qui est} é exclusivamente proprio de Deus, porque só em Deus o ser é a mesma essencia, e o nome, como dissemos, exprime a essencia. Ouçamos de novo S. Thomaz: « Ipsum esse Dei distinguitur et individuatur a quolibet alio esse per hoc ipsum, quod est esse per se subsistens, et non adveniens alicui naturae, quae sit aliud ab ipso esse. Omne autem aliud esse, quod non est subsistens, oportet quod individuetur per naturam et substantiam, quae in tali esse subsistit. Et in eis verum est, quod esse hujus est aliud ab esse illius per hoc quod est alterius naturae » (\emph{De pot}., q. 7, a. 2 ad 5).}.

26. \textbf{Especies de attributos divinos, derivados da essencia}. --- Os attributos divinos, derivados da essencia, dividem-se em \emph{quiescentes} e \emph{operativos}. --- Os \emph{quiescentes} são os que se referem á \pg{448} \pg{449}  \emph{essencia,} emquanto é uma \emph{entidade}. Os \emph{operativos} são os que se referem á \emph{essencia,} emquanto é \emph{principio de operações}, e comprehendem as \emph{faculdades}, com as \emph{operações} respectivas.

*

Todos os philosophos concordam em que em \textsc{Deus}, considerado segundo o nosso modo de entender, existe um attributo constitutivo da essencia divina, mas divergem na determinação d'aquelle attributo.

Ha quatro outras opiniões a este respeito.

\emph{a}) Uns dizem que o attributo \emph{constitutivo} da essencia divina é a somma \emph{actual} de todas as perfeições possiveis.

\emph{b}) Outros, como \emph{Scoto}, sustentam que a essencia divina não é constituida pela somma \emph{actual}, mas pela \emph{exigencia radical} de todas as perfeições.

\emph{c}) Outros, seguindo \emph{Billuart}, affirmam que o \emph{constitutivo} da essencia divina é a \emph{intelligencia}.

\emph{d}) Outros finalmente ensinam que o constitutivo da essencia de \textsc{Deus} é a \emph{aseidade} (\emph{ens a se}), isto é, aquella propriedade que faz com que \textsc{Deus} exista \emph{por si}, sem receber de outra causa o ser que possue.

\pg{450} Os attributos \emph{quiescentes} dividem-se em \emph{negativos} e \emph{positivos}. Os \emph{negativos} são os que \emph{negam} de \textsc{Deus} as imperfeições das creaturas. Os \emph{positivos} são os que \emph{affirmam} de \textsc{Deus} as perfeições creadas, mas isentas de todo o limite.

Neste artigo e no seguinte tratamos dos attributos \emph{quiescentes}, \emph{negativos} e \emph{positivos}; nos capitulos seguintes trataremos dos attributos \emph{operativos} \footnote{A nossa intelligencia, não podendo, por ser finita, comprehender a Essencia divina, que excede, pela sua perfeição infinita, a capacidade de toda a creatura, procura exprimil-a, de algum modo, por muitos e differentes conceitos. A esta pluralidade e diversidade de conceitos corresponde, da parte de Deus, não uma pluralidade ou diversidade real de perfeições, mas uma plena e infinita perfeição, que equivale a todas as perfeições e que as contém de um modo eminente. --- A esta pluralidade e diversidade de conceitos, que exprimem a unica e simplicissima Essencia divina, correspondem varios nomes, que por isso não são \emph{synonymos}, porque, embora todos exprimam a mesma Essencia, comtudo não a exprimem todos sob o mesmo aspecto, mas um exprime-a sob um aspecto e outro sob outro aspecto, e sempre de um modo \emph{deficiente} e \emph{parcial} (Cf. \emph{Sum}. \emph{Th}., p. I, q. 13, a. 4).}.

Todas estas opiniões são inacceitaveis.

a) E' inacceitavel a \emph{primeira}. Por quanto, a \emph{infinidade} das perfeições constitue a essencia \emph{physica} de \textsc{Deus}, mas não a \emph{metaphysica}; e aqui procuramos o \emph{constitutivo} da essencia \emph{metaphysica} de \textsc{Deus}.

\emph{b}) E' inacceitavel a \emph{segunda}. Com effeito, aquella \emph{existencia radical}, longe de constituir a essencia de \textsc{Deus}, suppõe-na já constituida.

c) E' inacceitavel a \emph{terceira}. Na verdade, a \emph{intelligencia} não é a \emph{primeira} coisa, que conhecemos de \textsc{Deus}, nem o distingue de todos os seres creados; pois ha muitas creaturas dotadas de intelligencia. --- Além d'isso, não podemos deduzir da \emph{intelligencia} muitos attributos de \textsc{Deus}, como a \emph{eternidade}, a \emph{immensidade}, a \emph{immutabilidade}, etc.

\emph{d}) E' inacceitavel a \emph{quarta}. Porquanto, a \emph{aseidade}, embora distinga Deus das creaturas, comtudo não representa o intrinseco, proximo e formal principio, que se concebe como constitutivo da essencia divina, e pelo qual Deus se distingue de todas as outras coisas. Na verdade, a independencia (\emph{esse} a se), embora exclua a causa efficiente, todavia não se concebe como uma coisa \emph{intrinseca} ao proprio ente, nem como a raiz de todos os attributos divinos, porque podemos sempre perguntar: porque Deus é \emph{ens a se}? E a resposta deve ser esta: porque \textsc{Deus} é o ser subsistente por si, e o acto purissimo. Por isso, este attributo, embora não seja o primeiro na ordem do \emph{conhecimento}, é o primeiro na ordem da \emph{realidade}, porque é o que se concebe como uma coisa \emph{intrinseca} e como a raiz e a fonte de todos os outros attributos divinos. --- Notamos que alguns philosophos e theologos tomam estas duas expressões --- \emph{Deus est ens per se subsistens} e \emph{Deus est ens a se} --- como equivalentes, e, dizendo que o constitutivo da essencia divina é a \emph{aseidade}, entendem que é a \emph{subsistencia do ser divino}. Mas, rigorosamente fallando, ha uma differença entre as duas expressões.

%%>>> art-40

\section*{Artigo II.\\ Attributos de Deus negativos.}
\addcontentsline{toc}{section}{Artigo II. --- Attributos de Deus negativos.}

\pg{451} \textbf{27. Attributos de Deus negativos.} --- Attributos \emph{negativos} são, como dissemos, os que \emph{negam} de \textsc{Deus} as imperfeições proprias das creaturas. --- Os principaes são: a \emph{simplicidade}, a \emph{immutabilidade}, a \emph{eternidade}, a \emph{immensidade} e a \emph{unidade} \footnote{Estes attributos são \emph{negativos} em quanto ao \emph{termo}, mas não em quanto á \emph{significação}, porque significam uma perfeição pura e simplicissima, de que a essencia de \textsc{Deus} é dotada. --- A \emph{infinidade} tambem colloca-se entre os attributos \emph{negativos}; mas d'ella fallámos no artigo precedente.}.

\textbf{28. Simplicidade de Deus.} --- A \emph{simplicidade} é o attributo, que exclue de \textsc{Deus} toda a \emph{composição} e \emph{divisibilidade} \footnote{\emph{Simplicidade}, em geral, é a \emph{negação de partes}; porque todos chamam \emph{simples} o ente, que não tem partes. A \emph{simplicidade} é uma especie de \emph{unidade}, á qual accrescenta a \emph{indivisibilidade}. \emph{Uno} é o ente, que é \emph{indiviso}. Mas o ente \emph{uno} pode ter partes, e, pela divisão, tornar-se multiplice; por isso, se exclue a multiplicidade \emph{actual}, não exclue a multiplicidade \emph{potencial}. Pelo contrario, o \emph{simples}, por não ter partes, exclue a multiplicidade \emph{actual} e a \emph{potencial}. --- A \emph{simplicidade}, considerada em si, não representa nem uma perfeição, nem uma imperfeição. Ás vezes o ente \emph{simples} é mais perfeito que o \emph{composto}, ás vezes é menos perfeito. E' mais perfeito, quando representa um \emph{acto} ou uma \emph{forma} completa e subsistente, dotada de toda a perfeição, propria da sua especie; assim o Anjo é mais perfeito que o homem. E' menos perfeito, quando representa uma \emph{parte} de um todo; assim a alma humana é menos perfeita que o composto humano. Diz S. Thomaz: « Simplicitas \emph{per se} non est causa nobilitatis, sed perfectio; unde ubi perfecta bonitas in uno simplici invenitur, simplex est nobilius, quam compositum: quando autem e contrario simplex est imperfectum, tunc compositum est nobilius, quam simplex, sicut homo nobilior est terra » (4 Sent. \emph{Dist}. 11, q. 2, a. 1 ad 1).\par A' \emph{simplicidade} oppõe-se diametralmente a \emph{composição}. --- Esta pode ser \emph{real}, ou \emph{de razão}. E' \emph{real}, quando a coisa é por si composta de partes, independentemente da nossa reflexão. E' \emph{de razão}, quando a nossa intelligencia, pela sua limitação natural, exprime por varios e diversos conceitos uma coisa, que por si é simplicissima. A composição \emph{real} denota imperfeição na coisa, e por isso não pode admittir-se em \textsc{Deus}; a \emph{de razão} denota apenas imperfeição na nossa intelligencia, e é compativel com a simplicidade de \textsc{Deus} (Cf. \emph{C}. \emph{Gent}., 1. II, c. 14, --- \emph{De ver}., q. 5, a. 2 ad 2). --- A composição \emph{real} subdivide-se em \emph{physica}, \emph{metaphysica} e \emph{logica}. A \emph{physica} é propria do ente, que consta de diversos elementos \emph{physicos}, como são a \emph{materia} e a \emph{forma}. A \emph{metaphysica} é propria do ente, que consta de partes \emph{metaphysicas}, como são a \emph{potencia} e o \emph{acto}, a \emph{essencia} e o \emph{supposto}, a \emph{essencia} e a \emph{existencia}, a \emph{essencia} e os \emph{attributos}, a \emph{substancia} e o \emph{accidente}. A \emph{logica} é propria do ente, que consta de elementos \emph{logicos}, como são o \emph{genero} e a \emph{differença}. --- Ninguem se admire se reduzimos a composição \emph{logica á real}. A composição \emph{logica}, se é \emph{de razão} emquanto ao seu conceito \emph{formal}, porque\pgn{452}os elementos são productos da intelligencia, é, comtudo, \emph{real} emquanto ao \emph{fundamento}, pois presuppõe uma composição \emph{real} no ente composto de \emph{genero} e de \emph{differença}, visto que o \emph{genero} se funda na \emph{potencia}, e a \emph{differença}, no \emph{acto}. --- Provaremos que \textsc{Deus} é \emph{simplicissimo}, se provarmos que Elle exclue de si toda e qualquer especie de \emph{composição}.}.

\pg{452} 29. \textbf{Deus} é \textbf{simplicissimo}.

a) Todo o ente composto é \emph{finito}; porque resulta de partes, e estas são essencialmente finitas. Ora \textsc{Deus} é \emph{infinito}. Logo \textsc{Deus} não é \emph{composto}, mas \emph{simplicissimo}.

b) O ente \emph{composto} não só é \emph{posterior}, ao menos por \emph{natureza}, ás partes componentes, mas tambem depende de uma causa \emph{extrinseca}, que reune as mesmas partes componentes. Ora \textsc{Deus} é o ente \emph{primeiro}, a causa independente de toda a causa. Logo \textsc{Deus} não é \emph{composto}, mas \emph{simplicissimo} \footnote{Cf. S. Thom., \emph{Sum Th.}, p. I, q. 3, a. 7; --- c. \emph{Gent.}, l. I, c. 16 e seg., --- \emph{De pot.} q. 7, a. 1. --- Depois de termos demonstrado em geral que em \textsc{Deus} não ha composição alguma, demonstraremos em particular que em \textsc{Deus} não ha composição \emph{physica}, nem \emph{metaphysica}, nem \emph{logica}.\par a) Em \textsc{Deus} não ha composição \emph{physica}. --- Em \textsc{Deus} haveria composição \emph{physica}, se elle fosse \emph{corpo}, ou \emph{forma corporea}, ou \emph{composto} de \emph{materia} e \emph{forma}. Ora \textsc{Deus} não é \emph{corpo}, nem \emph{forma} de corpo, nem \emph{composto}.\par \emph{a}) \textsc{Deus} não é \emph{corpo}. O corpo não é a mais nobre entre as substancias. Porquanto, o corpo \emph{vivo} é mais nobre do que o \emph{não-vivo}. Ora o corpo \emph{vivo} é mais nobre do que o \emph{não-vivo}, não emquanto é \emph{corpo} (aliás todos os corpos seriam eguaes em \emph{nobreza}), mas emquanto é \emph{vivo}. Logo o principio, de que deriva a \emph{vida} para o corpo, é mais nobre do que o \emph{corpo}, como tal. Ora \textsc{Deus} é o ser nobilissimo. Logo \textsc{Deus} não é \emph{corpo} (\emph{Sum. Th.}, p. I, q. 3, a. 1).\par \emph{b}) \textsc{Deus} não é \emph{forma} nem \emph{materia} de algum corpo. --- Não é \emph{forma}. Porquanto, a \emph{forma corporea}, além de ser \emph{finita} e por isso \emph{creada} e \emph{contingente}, depende, pelo menos \emph{extrinsecamente} (como a alma humana), da materia. Ora \textsc{Deus} é infinito, Causa primeira, Ente necessario e absolutamente independente. Logo não pode ser \emph{forma} de algum corpo. --- Muito menos pode ser \emph{materia}; pois esta é um ser \emph{potencial}, \emph{finito}, \emph{contingente}, e por isso dotada de propriedades contrarias ás de Deus (\emph{Sum. Th.}, p. I, q. 3 a. 8).\par c) \textsc{Deus} não é \emph{composto} de \emph{materia} e \emph{forma}. Com effeito, todo o \emph{composto}, além de ser \emph{finito}, \emph{é dependente} não só da causa \emph{efficiente}, que reune as partes, mas tambem da \emph{forma}, de que recebe as perfeições. Ora Deus não depende de nenhuma causa, por ser a Causa primeira, e possue todas as perfeições não por \emph{participação}, mas por \emph{essencia}. (\emph{Sum. Th.}, p. I, q. 3, a. 2). --- E', pois, absurdo o \emph{anthropomorphismo}, que attribue a \textsc{Deus} a \emph{forma humana}.\par D'isto resulta que \textsc{Deus} é \emph{purissimo espirito}. \emph{Espirito} é uma substancia independente da \emph{materia}. Ora a substancia divina exclue toda e qualquer mistura de \emph{materia} e de \emph{potencia}. Logo \textsc{Deus} é \emph{purissimo espirito}.\par Totavia, do facto de \textsc{Deus} ser \emph{purissimo espirito}, e não ter \emph{corpo}, não póle concluir-se (como concluem os \emph{pantheistas}) que Elle não é \emph{pessoa}. Por quanto, \emph{pessoa} é a \emph{substancia individua de natureza racional}. Ora \textsc{Deus} é \emph{substancia}, subsistente em si e por si, e esta \emph{substancia} é infinitamente \emph{intelligente}. Logo \textsc{Deus} é \emph{pessoa}. --- E não só é \emph{pessoa}, mas possue a \emph{persona}\pgn{453}\emph{lidade} no grau mais elevado, que póde imaginar-se; porque a \emph{substancia} de \textsc{Deus} é \emph{perfeitissima}, e a sua intelligencia é \emph{infinita}. --- A materia não é sempre elemento necessario da \emph{personalidade}; pelo contrario, sendo a \emph{razão} o elemento principal da \emph{pessoa}, é claro que, quanto mais a razão for independente da materia, tanto mais nobre será a \emph{pessoa}.\par b) Em \textsc{Deus} não ha composição \emph{metaphysica}. --- Na verdade, \textsc{Deus} não é composto de \emph{potencia} e \emph{acto}, de \emph{essencia} e \emph{supposto}, de \emph{essencia} e \emph{existencia}, de \emph{essencia} e \emph{attributos}, de \emph{substancia} e \emph{accidentes}.\par \emph{a}) \textsc{Deus} não é composto de \emph{potencia} e \emph{acto}. Todo o ente, composto de \emph{potencia} e \emph{acto}, \emph{é mutavel}. Ora \textsc{Deus} é \emph{immutavel}. Logo não é composto de \emph{potencia} e \emph{acto}. --- Por isso \textsc{Deus} é \emph{acto purissimo}, contendo \emph{todas} as \emph{possiveis} perfeições num grau \emph{infinito}; porque, se não possuisse alguma perfeição ou fosse capaz de augmento, estaria em \emph{potencia} para aquella perfeição ou augmento de perfeição, e por isso seria \emph{composto} de \emph{potencia} e \emph{acto}.\par \emph{b}) \textsc{Deus} não é composto de \emph{essencia} e de \emph{supposto}. Um ente é composto de \emph{essencia} e de \emph{supposto}, quando os elementos, que estão contidos neste, differem, na \emph{especie} ou no \emph{numero}, dos que estão contidos naquella; de modo que não se póde dizer que o \emph{supposto} é a \emph{essencia}. Assim o homem é composto de \emph{essencia} e de \emph{supposto}; porque, além de \emph{todos} os elementos \emph{especificos}, que estão contidos na \emph{humanidade} e que são a \emph{animalidade} e a \emph{racionalidade}, o \emph{supposto humano} contém outros elementos, que são as \emph{propriedades individuaes} e que não estão contidas no conceito \emph{especifico} de \emph{humanidade}; de modo que não podemos dizer que \emph{Pedro} é a \emph{humanidade}, mas só que \emph{Pedro} possue a \emph{humanidade}. Ora as \emph{propriedades individuaes} derivam da \emph{materia}. Logo nos \emph{espiritos puros}, em que não ha \emph{materia}, tudo o que está no \emph{supposto}, encontra-se na \emph{essencia}; de modo que esta póde attribuir-se áquelle por \emph{identidade}, dizendo que o \emph{supposto} é a propria \emph{essencia}. \textsc{Deus} é \emph{purissimo espirito}. Logo, em \textsc{Deus}, não ha distincção entre \emph{supposto} e \emph{essencia}, e por isso \textsc{Deus} não é composto de \emph{supposto} e \emph{essencia}; de modo que podemos dizer que \textsc{Deus} é a \emph{Divindade} (\emph{Sum. Th.}, p. I, q. 3, a. 3; --- c. \emph{Gent.} l. I, c. 21).\par \emph{c}) \textsc{Deus} não é composto de \emph{essencia} e de \emph{existencia}. Um ente, composto de \emph{essencia} e de \emph{existencia}, isto é, cuja \emph{essencia} se distingue \emph{realmente} da \emph{existencia}, não existe \emph{essencialmente}, por assim o exigir a sua essencia, mas recebeu o \emph{ser} de outra causa. Ora \textsc{Deus} existe \emph{essencialmente}, não recebendo de outra causa o ser que possue. Logo \textsc{Deus} não é composto de \emph{essencia} e de \emph{existencia}. --- Além disso, a \emph{essencia} refere-se á \emph{existencia}, como a \emph{potencia} ao \emph{acto}. Ora \textsc{Deus} não é composto de \emph{potencia} e \emph{acto} (\emph{Sum. Th.}, p. I, q. 3, a. 4; --- c. \emph{Gentes}, l. I, c. 22).\par \emph{d}) \textsc{Deus} não é composto de \emph{essencia} e \emph{attributos}. A Essencia divina, sendo o proprio ser, possue a plenitude da perfeição, e por isso não ha attributo algum, que Ella não contenha. Ora a \emph{essencia} de \textsc{Deus}, se se distinguisse \emph{realmente} dos seus attributos, aperfeiçoar-se-hia pelos mesmos attributos, e já não possuiria a plenitude da perfeição. Logo em \textsc{Deus} não ha distincção \emph{real} entre a \emph{essencia} e os \emph{attributos}. Logo \textsc{Deus} não é composto de \emph{essencia} e de \emph{attributos}. --- Além disso, quando os attributos são \emph{realmente} distinctos da essencia, esta refere-se áquelles, como a potencia ao acto. Ora a Essencia divina, por ser o proprio ser subsistente, é acto purissimo, sem sombra de\pgn{454}potencialidade. Diz S. Thomaz: « Perfectio divini \emph{esse} non attenditur secundum aliquod \emph{additum} supra ipsum; sed quia ipsum secundum \emph{seipsum} perfectum est. Bonitas igitur Dei non est aliquid additum suae substantiae, \emph{sed} sua \emph{substantia est sua bonitas} » (\emph{C. Gentes}, l. I, c. 38). --- Se, em \textsc{Deus}, os attributos são uma e a mesma coisa com a essencia, é claro que, assim como dizemos que \textsc{Deus} é a \emph{Divindade}, assim tambem devemos dizer que Elle é a \emph{Verdade}, a \emph{Belleza}, a \emph{Bondade}, a \emph{Misericordia}, etc.\par Se, em Deus, os attributos não se distinguem \emph{realmente} da \emph{essencia}, tambem não se distinguem \emph{realmente} entre si, e por isso, na realidade, um identifica-se com outro. E' um corollario da these precedente. Se cada attributo divino se identifica com a Essencia, é evidente que um não se distingue realmente do outro, porque uma coisa não se distingue realmente de si mesma. A Essencia divina encerra, na sua absoluta simplicidade, todas as perfeições possiveis; e só a imperfeição e limitação da nossa intelligencia considera como derivadas da essencia e distinctas entre si as propriedades, que, na realidade, são uma e a mesma coisa. --- E' claro que isto se refere aos attributos \emph{absolutos} de Deus, e não aos que se dizem \emph{relativos ad intra} e que são oppostos entre si, como são o \emph{Paternidade}, a \emph{Filiação} e a \emph{Processão}; porque estes attributos são realmente distinctos e oppostos e constituem as tres Pessoas da Santissima Trindade --- \textsc{Padre}, \textsc{Filho} e \textsc{Espirito Santo}. Diz S. Thomaz: « Omnia in Deo sunt unum \emph{re}, praeter ingenerationem, generationem et processionem, quae constituunt personas re distinctas » (I \emph{Sent}. \emph{Dist}. 2, q. 1, a. 3 ad 6).\par Mas, se os attributos divinos não se distinguem \emph{realmente} um do outro, porque cada um se identifica com a Essencia divina, comtudo distinguem-se \emph{logicamente}, porque o conceito de um não se identifica com o conceito de outro; assim o conceito de \emph{justiça} não é o conceito de \emph{misericordia}. E esta distincção \emph{logica} não é arbitraria, mas tem o seu \emph{fundamento real} na propria Essencia divina. Porquanto, a cada um d'esses diversos conceitos, relativos a attributos diversos, corresponde, na sua plenitude simplicissima, a unica Essencia divina. Se assim não fosse, isto é, se de nenhum modo existisse em Deus o fundamento d'essa distincção, os nomes \emph{absolutos} e \emph{positivos}, que exprimem as divinas perfeições, seriam \emph{synonymos}; o que é falso; pois esses nomes, embora exprimam todos a Essencia divina, exprimem-na comtudo sob differentes aspectos; assim a \emph{justiça} exprime a Essencia divina emquanto dá a cada um o que merece, a \emph{misericordia} exprime a mesma Essencia divina emquanto tem dó dos nossos males. E' a doutrina de S. Thomaz. Diz assim: « In Deo est sapientia, bonitas, et hujusmodi, quorum quodlibet est ipsa divina essentia, et ita sunt unum re. Et quia unumquodque eorum est in Deo secundum verissimam sui rationem, et ratio sapientiae non est ratio bonitatis in quantum hujusmodi, relinquitur quod sunt diversa ratione non tantum ex parte ipsius ratiocinantis, sed ex proprietate ipsius rei; et inde est quod Ipse non est causa rerum omnino aequivoca, cum secundum formam suam producat effectus similes non univoce sed analogice » (I \emph{Sent}. \emph{Dist}. 2, q. 1, a. 2; cf. \emph{Sum. Th.}, p. I, q. 13, a. 2). E mais adeante: « Pluralitas nominum non tantum est ex parte intellectus nostri, formantis diversas conceptiones de Deo; sed etiam ex parte ipsius Dei, inquantum scilicet est\pgn{455}aliquid in Deo correspondens omnibus istis conceptionibus, scilicet plena et omnimoda ejus perfectio » (I \emph{Sent}. \emph{Dist}. 2, q. 1, a. 3).\par e) \textsc{Deus} não é composto de \emph{substancia} e \emph{accidentes}. O \emph{accidente} refere-se á \emph{substancia}, como o \emph{acto} á \emph{potencia}. Ora \textsc{Deus} é \emph{acto purissimo}. --- Além d'isso, o \emph{accidente} aperfeiçoa a \emph{substancia}. Ora a \emph{substancia divina}, possuindo a plenitude da perfeição, não pode receber novas perfeições. Logo \textsc{Deus} não é composto de \emph{substancia} e de \emph{accidentes}.\par c) Em \textsc{Deus} não existe composição \emph{logica}, isto é, de \emph{genero} e \emph{differença}. Por quanto o \emph{genero} refere-se á \emph{differença}, como a \emph{potencia} ao \emph{acto}. Logo \textsc{Deus}, por ser \emph{acto purissimo}, não é composto de \emph{genero} e de \emph{differença}. --- Além d'isso, o ente, composto de \emph{genero} e de \emph{differença}, é \emph{finito}, porque só possue a perfeição propria de uma certa especie de entes. Ora \textsc{Deus} é \emph{infinito}. --- Mais, a composição de \emph{genero} e de \emph{differença}, embora seja \emph{intencional}, comtudo suppõe uma composição \emph{real} de \emph{potencia} e de \emph{acto}; pois o \emph{genero} funda-se na \emph{potencia}, e a \emph{differença} no \emph{acto}. Ora em \textsc{Deus} não ha composição \emph{real}, e por isso não pode haver composição \emph{logica}. --- Se \textsc{Deus} podesse estar subordinado a um \emph{genero}, este seria o \emph{ente} ou a \emph{substancia}. Ora não pode ser o \emph{ente}; porque este não se attribue a \textsc{Deus} e ás creaturas em significação \emph{univoca}; visto que, quando se attribue a \textsc{Deus}, \emph{ente} significa o proprio ser subsistente em si, e, quando se attribue ás creaturas, significa o ser recebido num sujeito ou potencia, que é a essencia. Nem pode ser a \emph{substancia}; porque a \emph{substancia}, que é \emph{genero}, não é o proprio ser, aliás toda a substancia seria o proprio ser, isto é, o seu ser não seria produzido por outra causa; e por isso \textsc{Deus}, que é o proprio ser, não pode estar sujeito ao \emph{genero} de \emph{substancia}. --- Todavia, \emph{analogicamente} e num sentido lato, podemos dizer que \textsc{Deus} é \emph{substancia}, emquanto esta significa um ente \emph{subsistente}. --- D'ahi se vê que não podemos formar uma \emph{rigorosa} definição de \textsc{Deus}, porque uma definição \emph{rigorosa} é composta de \emph{genero} e de \emph{differença}. Mas nem por isso confundimos \textsc{Deus} com as creaturas. \textsc{Deus} distingue-se das creaturas --- não por differenças especificas e componentes, --- mas pela pureza absoluta, pela substancia simplicissima do proprio ser (\emph{Sum. Th.}, p. I, q. 3, a. 5).}.

\pg{453} 30. \textbf{Immutabilidade de Deus}. --- \emph{Immutabilidade} é o attributo, que afasta de \textsc{Deus} toda a \emph{mudança}. --- A \emph{mudança} pode ser

\pg{454} duplice: \emph{physica} ou \emph{moral}. A mudança \emph{physica} consiste na adquisição ou na perda de alguma entidade ou perfeição; a

\pg{455} \emph{moral} consiste na passagem da \emph{vontade} de um para outro acto,

de um para outro proposito \footnote{S. Thomaz trata d'isto --- na \emph{Sum. Th.}, p. I, q. 9; --- q. 14, a. 15; q. 19, a. 7; --- \emph{c. Gentes}, l. I, c. 13 e seg.; --- I \emph{Sent.} Dist. 8, q. 3, a. 1 e 2.}.

31. \textbf{Deus é immutavel.}

a) \textsc{Deus} é \emph{absolutamente} immutavel. Todo o ente, que se muda de algum modo, é composto de \emph{potencia} e de \emph{acto}. Ora \textsc{Deus} é acto simplicissimo e purissimo, sem sombra de potencialidade. Logo \textsc{Deus} é \emph{absolutamente} immutavel.

b) \textsc{Deus} é \emph{physicamente} immutavel. Todo o ente, que se muda \emph{physicamente}, ou adquire uma perfeição que não tinha, ou perde uma perfeição que tinha. Ora \textsc{Deus}, sendo o Ente \emph{infinito} e \emph{necessario}, possue todas as perfeições possiveis, e não póde perdel-as, aliás perderia a essencia, com que as perfeições se identificam. Logo \textsc{Deus} é \emph{physicamente} immutavel.

\pg{456} c) \textsc{Deus} é \emph{moralmente} immutavel. Com effeito, a mudança \emph{moral}, ou de \emph{vontade}, deriva de tres causas, --- ou da mudança \emph{physica} do sujeito, que, por isso, começa a desejar uma coisa, pela qual tinha repugnancia, --- ou da mudança do \emph{conhecimento}, emquanto a intelligencia descobre uma causa, que antes ignorava, e assim determina a vontade a mudar de proposito, --- ou da leviandade e inconstancia da propria \emph{vontade}, que, sem motivo sufficiente, passa de um proposito para outro. Ora nenhuma d'estas tres causas pode admittir-se em \textsc{Deus}, que, sendo infinitamente perfeito, não está sujeito a mudanças \emph{physicas}, nem á ignorancia ou inconstancia. Logo \textsc{Deus} é \emph{moralmente} immutavel \footnote{D'onde resulta que \textsc{Deus}, quando quer ou permitte effeitos diversos ou contrarios entre si, não muda de vontade; porque uma coisa é mudar de vontade e outra coisa é querer a mudança no effeito, e um agente póde, com o mesmo decreto, prescrever que presentemente se faça uma coisa, e que depois se faça o contrario. --- Como tambem, quando criou o mundo, \textsc{Deus} não mudou de vontade, porque não mudou de essencia (visto que em \textsc{Deus}, como provaremos, as \emph{operações} se identificam com a \emph{essencia}), mas, desde a eternidade, decretou que, em um certo tempo, existisse o mundo. Diz S. Thomaz: « Aliud est mutare voluntatem, et aliud est velle aliquarum rerum mutationem. Potest enim aliquis, eádem voluntate immobiliter permanente, velle quod nunc fiat hoc, et postea fiat contrarium. Sed tunc voluntas mutaretur, si aliquis inciperet velle quod prius non voluit. Quod quidem accidere non potest, nisi praesupposita mutatione vel ex parte cognitionis, vel circa dispositionem substantiae ipsius volentis » (\emph{Sum. Th.}, p. I, q. 19, a. 7). --- Entende-se egualmente como é que \textsc{Deus}, sem alguma mudança na sua vontade, possa agora amar e amanhã deixar de amar uma e a mesma pessoa, que hoje é santa e amanhã é peccadora; porque a mudança não se deu na vontade de \textsc{Deus}, mas só no homem, que antes era digno de amor, e depois se tornou digno de odio. Diz S. Agostinho: « Cum iratus malis dicitur (Deus), et placidus bonis, illi mutantur, non Ipse: sicut lux infirmis oculis aspera, firmis lenis est, ipsorum scilicet mutatione, non sua » (\emph{De Trin.}, l. V, c. 16).\par Objectam: \textsc{Deus} é livre com relação ás creaturas. A liberdade suppõe a indifferença; de modo que o acto, que se diz livre, pode ser e pode não ser realizado. Ora se \textsc{Deus} pode realizar e pode deixar de realizar um acto, é evidente que, quando o realiza, a sua vontade se muda.\par Resposta. Tudo o que \textsc{Deus} quer fóra da sua essencia, não só o quer desde a eternidade, mas tambem com o mesmo simplicissimo acto, com que ama a si mesmo; pois, em \textsc{Deus}, não ha nem pode haver composição ou vicissitude de actos. O acto da divina vontade, sendo simplicissimo, refere-se ao objecto de um modo proporcionado á condição do mesmo objecto. Logo \textsc{Deus} ama a si mesmo como \emph{fim}, e ama as outras coisas como \emph{meios} para o mesmo fim. E, como estes meios não têm uma relação necessaria com o fim, \textsc{Deus} pode querel-os e deixar de querel-os; e, ainda que ame a si mesmo \emph{necessariamente}, comtudo ama ou quer as outras coisas \emph{contingente} e \emph{livremente}, porque assim o exige a diversa qualidade dos objectos. --- Por isso\pgn{457}Deus quer verdadeiramente as creaturas, quando as quer; e não as quer verdadeiramente, quando não as quer; mas o querer e o não querer de Deus não suppõe uma mudança da parte de Deus, mas da parte das creaturas, que se referem \emph{realmente} ao proprio Deus. --- Para a comprehensão d'esta verdade, lembramos que a causa efficiente creada não se muda, porque opera; mas muda-se, porque, operando, passa necessariamente da potencia para o acto; e por isso Deus, que é acto purissimo, não se muda quando opera. --- Em poucas palavras: a liberdade, no seu conceito abstractissimo, exprime indifferença não para com a propria \emph{entidade}, mas para com o \emph{termo}, de modo que a vontade possa querer ou deixar de querer um objecto; mas não exige que este querer ou deixar de querer seja expressão de actos \emph{entitativamente} diversos, se não quando ella se considera na vontade creada, a qual, pela sua volição, passa da potencia para o acto.\par \emph{Renan}, seguido pelos defensores do \emph{progresso indefinido}, distingue um duplice \textsc{Deus}, --- o \textsc{Deus} \emph{ideal} e o \textsc{Deus} \emph{real}. O \textsc{Deus} \emph{ideal} é, segundo esse escriptor, immutavel e eterno; o \textsc{Deus} \emph{real}, que se identifica com a \emph{universalidade das coisas existentes}, \emph{é mutavel} e capaz de progresso indefinido.\par Estas impias e absurdas affirmações, que encerram o mais crasso \emph{pantheismo}, não merecem ser novamente refutadas, pois já o foram no capitulo antecedente, quando provámos que existe um \textsc{Deus} \emph{real} e \emph{verdadeiro}.}.

\pg{457} 32. \textbf{Eternidade de Deus}. --- A \emph{eternidade} é o attributo, que exclue de \textsc{Deus} o \emph{principio}, o \emph{fim}, e toda a \emph{successão} \footnote{\emph{Boecio} dá a seguinte definição da \emph{eternidade}: « A \emph{eternidade é a posse total e perfeita de uma vida sem fim} ». Analysemos esta definição. A \emph{eternidade} --- é \emph{posse}, porque é uma coisa estavel, que se conserva com socego e não póde ser tirada por nenhuma força extranha; --- é posse \emph{total}, porque a duração de \textsc{Deus} é sempre completa e indivisivel, e distingue-se da duração do ser finito, a qual não é \emph{total} mas \emph{successiva}; --- \emph{perfeita}, porque a \emph{eternidade} é o \emph{presente}, que dura sempre, e por isso é diversa do tempo \emph{presente} dos seres finitos, que, embora, por não ter partes, seja \emph{total}, comtudo não é uma coisa absolutamente perfeita (pois não é perfeição possuir uma coisa apenas por um instante, e o tempo \emph{presente} das coisas finitas é um instante); --- é posse da uma \emph{vida}; porque a duração refere-se principalmente ao \emph{acto} dos seres nobilissimos, que é \emph{viver}. --- de uma vida \emph{sem fim}, isto é, que não principiou nem ha de acabar. --- A \emph{posse} sem \emph{principio} e sem \emph{fim} é o elemento \emph{material} da \emph{eternidade}; a \emph{simultaneidade} da posse é o elemento \emph{formal}.\par \emph{A eternidade}, pois, consiste na exclusão do \emph{principio}, do \emph{fim}, e de toda a \emph{successão}, e identifica-se com a propria substancia de \textsc{Deus}, emquanto Esta não teve \emph{principio}, nem terá \emph{fim}, e não está sujeita a \emph{successão}.\par Por isso o \emph{tempo} não é uma parte da \emph{eternidade}, porque a \emph{eternidade} não tem \emph{partes}; mas a \emph{eternidade} corresponde e assiste a todos os seculos, que se vão succedendo, assim como um \emph{centro} immovel corresponde, por meio de linhas, a todos os pontos de uma esphera, que se mova em volta d'elle.\par D'onde se vê o erro --- tanto de \emph{Locke}, quando disse que a \emph{eternidade} é constituida pela infinita multiplicação do \emph{tempo}, --- como de \emph{Clarke}, quando opinou que a \emph{eternidade} é o proprio tempo, considerado em \emph{abstracto}.\par S. Thomaz trata d'este argumento --- na \emph{Summ. Theol.}, p. I, q. 10; --- \emph{c. Gentes}, l. I, c. 15; --- \emph{I Sent.} Dist. 8, q. 2, a. 1, e Dist. 19, q. 2, a. 1 e 2.}

\pg{458} 33. \textbf{Deus} é \textbf{eterno}. --- Um ente e é diz-se \emph{eterno}, quando não tem \emph{principio}, nem \emph{fim}, e não está sujeito a \emph{successão}. Ora \textsc{Deus} não tem \emph{principio} nem \emph{fim}, e não está sujeito a \emph{successão}. --- Não tem \emph{principio} nem \emph{fim}. Na verdade, \textsc{Deus}, sendo o Ente \emph{necessario}, \emph{immutavel} e \emph{independente}, existiu sempre e não póde deixar de existir; pois, se tivesse principiado a existir, quando não existia, ou podesse deixar de existir, seria \emph{contingente}, \emph{mutavel} e \emph{dependente}. --- Não está sujeito a \emph{successão}. Porquanto, a \emph{successão} indica \emph{mudança} e \emph{imperfeição}, e \textsc{Deus} é \emph{immutavel} e infinitamente \emph{perfeito}. Logo \textsc{Deus} é \emph{eterno} \footnote{Dizemos no texto que toda a \emph{successão} indica \emph{imperfeição}; porquanto o ente, que admitte \emph{successão} nos seus estados ou qualidades, adquire perfeições que não tinha, ou perde as que tinha, e por consequencia é finito ou limitado nas suas perfeições.\par \textsc{Deus} não só é \emph{eterno}, mas é a sua \emph{eternidade}; porque a eternidade não é uma coisa distincta de Deus, mas é o seu proprio ser uniforme; e por isso Deus, assim como é a sua essencia, assim tambem é a sua eternidade. Diz S. Thomaz: « Ratio aeternitatis consequitur immutabilitatem, sicut ratio temporis consequitur motum. Unde cum Deus sit maxime immutabilis, sibi maxime competit esse aeternum. Nec solum est aeternus, sed est sua aeternitas; cum tamen nulla alia res sit sua duratio, quia non est suum \emph{esse}. Deus autem est suum esse uniforme. Unde: sicut est sua essentia, ita sua aeternitas » (\emph{Sum}. \emph{Th}., p. I, q. 10, a. 2).\par Como a \emph{eternidade} consiste na perfeita immutabilidade, e todas as creaturas são mutaveis, segue-se que a nenhuma creatura compete a eternidade, no sentido rigoroso da palavra, e por isso só \textsc{Deus} é eterno. --- Todavia costumamos dizer que algumas creaturas participam da eternidade, como participam da immutabilidade no seu ser ou nos seus actos. Mas trata-se sempre de uma eternidade impropria e imperfeita (\emph{Sum}. \emph{Th}., p. I, q. 10, a. 3).}.

34. \textbf{Immensidade de Deus}. --- \emph{Immensidade} é o attributo, pelo qual \textsc{Deus} está presente \emph{actualmente} em todas as coisas \emph{existentes} e \emph{virtualmente} em todas as coisas \emph{possiveis} \footnote{Quando dizemos que \textsc{Deus} está \emph{virtualmente} presente em todas as coisas \emph{possiveis}, entendemos que \textsc{Deus} tem a \emph{virtude} ou o poder de estar em todas as coisas, que não existem, mas que podem existir.\par E' necessario distinguir a \emph{immensidade} da \emph{ubiquidade}. --- A \emph{immensidade} é a presença divina, \emph{actual} ou \emph{virtual}, em todas as coisas, tanto \emph{existentes} como \emph{possiveis}; a \emph{ubiquidade} é a \emph{actual} presença de \textsc{Deus} nas coisas \emph{existentes}. --- A \emph{immensidade} é uma coisa eterna, absoluta e infinita, porque se extende a todo o possivel, e prescinde da \emph{actual} existencia das coisas; a \emph{ubiquidade} é uma coisa temporanea, relativa, não é rigorosamente infinita, porque suppõe a existencia do mundo, e porque as coisas creadas não são infinitas. --- A \emph{ubiquidade} está contida na \emph{immensidade}. --- Por isso \textsc{Deus},\pgn{459}embora não esteja presente nas coisas finitas, senão depois da creação, todavia é \emph{immenso} desde a eternidade; porque a propria Essencia de \textsc{Deus} exige que Elle possa e deva estar presente em todas as coisas, \emph{existentes} e \emph{possiveis}. --- S. Thomaz trata d'este assumpto --- na \emph{Sum}. \emph{Th}., p. I, q. 8, --- c. \emph{Gentes}, l. III, c. 68; --- I \emph{Sent}. Dist. 37, q. 3, a. 1.}.

35. \textbf{Deus} é immenso. --- \textsc{Deus} será \emph{immenso}, se estiver presente --- não só em todas as coisas \emph{creadas} e em todos os logares \emph{existentes}, --- mas tambem em todas as coisas e logares \pg{459}  \emph{possiveis}. Ora \textsc{Deus} --- está presente em todas as coisas \emph{creadas}, porque Elle produziu e conserva todas as creaturas, e um agente não póde operar onde não existe; --- está presente em todos os logares \emph{existentes}, porque todo o logar é uma coisa creada; --- póde estar presente em todas as coisas e logares \emph{possiveis}, porque \textsc{Deus} deve estar em todo o ser, a que se extende o seu poder (pois, em \textsc{Deus}, os attributos identificam-se com a substancia), e o poder de \textsc{Deus} extende-se a todo o \emph{possivel}. Logo \textsc{Deus} é \emph{immenso} \footnote{\textsc{Deus} está em todos os logares --- com o seu \emph{poder}, emquanto tudo Lhe obedece, --- com a sua \emph{presença}, emquanto Elle vê tudo o que se faz e pensa, --- com a sua \emph{essencia}, emquanto produz e conserva o ser das creaturas (Cf. \emph{Sum}. \emph{Th}., p. I, q. 8, a. 3).\par Pergunta-se: Onde estava \textsc{Deus}, antes da creação do mundo? e agora como está Elle no mundo?\par Antes da creação, \textsc{Deus} estava em si mesmo. E, quando o mundo foi creado, \textsc{Deus} não deixou de estar em si mesmo, nem se mudou pelo facto de estar presente em todas as coisas; assim como não se muda o sol, quando se mudam as coisas, que elle illumina. --- E, agora, \textsc{Deus} está no mundo, não diffundido por toda a parte, como a luz ou o ar, mas de um modo ineffavel, que excede toda a imaginação. Uma semelhança, embora infinitamente \emph{inadequada}, da presença de \textsc{Deus} nas coisas é a nossa alma, a qual, ainda que \emph{simples} e \emph{espiritual}, está toda em todo o organismo e toda em cada parte d'elle. Dizemos --- infinitamente \emph{inadequada}, porque a nossa alma é limitada ao nosso corpo, ao passo que \textsc{Deus} não é limitado a este mundo.\par Muitos negaram ou desvirtuaram a \emph{immensidade} de \textsc{Deus}.\par \emph{a}) Uns, como os \emph{Manicheus} e \emph{Descartes}, negaram a \emph{ubiquidade} de \textsc{Deus}; pois aquelles diziam que só as coisas \emph{espirituaes} estão sujeitas ao \emph{poder} de \textsc{Deus}, e \emph{Descartes} affirmou que \textsc{Deus} está em todas as coisas com o seu \emph{poder} e \emph{presença}, mas não com a sua \emph{essencia} ou \emph{substancia}.\par Ambas essas affirmações são falsas. --- E' falsa a dos \emph{Manicheus}; porque o poder de \textsc{Deus} é infinito e extende-se a todas as creaturas, espirituaes e corporeas. --- E' falsa a de \emph{Descartes}; porque, identificando-se em \textsc{Deus} as operações com a essencia, onde existem aquellas, deve existir esta.\par \emph{b}) Outros, seguindo as theorias de \emph{Newton} e \emph{Clarcke}, disseram que a \emph{immensidade divina} é o espaço infinito.\par Esta opinião foi refutada na \emph{Ontologia}.\par \emph{c}) \emph{Spinoza} disse que a \emph{immensidade} não é o attributo de um Ente individuo e pessoal, mas é a \emph{collecção} de todas as coisas.\par Estas affirmações contêm o \emph{pantheismo}, que já foi refutado.\par \emph{d}) \emph{Bayle} diz que \textsc{Deus}, sendo \emph{simplicissimo}, não póde ser \emph{immenso}; porque a \emph{simplicidade} e a \emph{immensidade} são attributos contradictorios.\par Respondemos que a \emph{simplicidade} não só não é contraria á \emph{immensidade},\pgn{460}mas é uma condição indispensavel para ella; porque \textsc{Deus} não seria \emph{immenso}, se não fosse \emph{simplicissimo}, pois o \emph{composto}, resultando de partes necessariamente finitas e limitadas, não póde ser \emph{immenso}.}.

\pg{460} 36. \textbf{Unidade de Deus}. --- \emph{Unidade} é o attributo, pelo qual a Essencia divina não póde dividir-se, multiplicar-se e constituir varios e diversos deuses \footnote{Quando, pois, dizemos que \textsc{Deus} é \emph{uno}, entendemos que é uno não só \emph{especifica} mas tambem \emph{numericamente}, e por isso é \emph{unico}.\par S. Thomaz trata d'este assumpto --- na \emph{Sum Th}., p. I, q. 11, a. 3 e 4; --- c. \emph{Gentes}, l. I, c. 42; --- \emph{De pot}., q. 3, a. 6; --- I \emph{Sent}. Dist. 2, a. 1, --- II \emph{Sent}. Dist. 21, q. 1, a. 1.}.

37. \textbf{Deus} é uno.

a) \textsc{Deus} é o Ente infinitamente perfeito. Ora o ente infinitamente perfeito deve ser \emph{uno}. Porquanto, se fossem varios, deveria haver entre elles uma distincção. Esta distincção só poderia dar-se, emquanto um fosse dotado de uma perfeição, que faltasse a outro. Mas, neste caso, aquelle, a quem faltasse uma perfeição, não seria \textsc{Deus}; mas só o seria Aquelle, que fosse infinitamente perfeito. Logo \textsc{Deus} é \emph{uno}.

b) Uma forma multiplica-se, emquanto é recebida em muitos e diversos sujeitos; assim a \emph{humanidade} multiplica-se, emquanto é recebida em muitos e diversos homens. Se, porém, a forma não for recebida em sujeito algum, mas subsistir em si mesma, não poderá multiplicar-se, e existirá um unico individuo d'aquella especie; assim, se a \emph{humanidade} fosse subsistente em si e por si, existiria um unico homem. Ora a Divindade não é recebida em sujeito algum (porque, em \textsc{Deus}, infinitamente \emph{simples}, a \emph{essencia} identifica-se com o \emph{supposto}), mas é subsistente em si e por si. Logo a Divindade não póde multiplicar-se e formar varios suppostos. Logo \textsc{Deus} é \emph{uno} \footnote{Nunca devemos esquecer que a \emph{subsistencia do ser divino} é, segundo o nosso modo de entender, a razão e a raiz de todos os attributos divinos. O ser das creaturas é simples como é simples uma parte ou uma coisa incompleta, porque é uma forma recebida num sujeito e individualisada pela natureza, que subsiste nesse ser, e por isso é multiplicavel e multiplica-se conforme a multidão das creaturas. Pelo contrario, o ser de Deus é simples, como é simples o acto e a perfeição, porque é o ser por si subsistente e por si individualisado e distincto de todo e qualquer outro ser, e por isso não é multiplicavel, mas é \emph{unico}.}.

c) Se existissem varios deuses, \emph{especificamente} eguaes, mas \emph{numericamente} distinctos, seriam necessariamente compostos de \emph{natureza divina}, \emph{commum} a todos, e de \emph{notas} ou \emph{attributos}, \emph{proprios} de cada um (aliás seria impossivel a egualdade \emph{especifica} a e distincção \emph{numerica}). Ora \textsc{Deus}, sendo \emph{simplicissimo}, exclue toda \pg{461}  e qualquer composição. Logo não podem existir varios deuses. Logo \textsc{Deus} é \emph{uno} \footnote{Os systemas ou crenças contrarias á \emph{unidade} de \textsc{Deus} foram o \emph{dualismo} e o \emph{polytheismo}.\par a) O \emph{dualismo} admittia que existem \emph{dois principios infinitos}, \emph{necessarios}, e \emph{supremos} --- um \emph{bom} e outro \emph{mau}, e que do \emph{principio bom} derivam todos os \emph{bens} e do \emph{mau} todos os \emph{males}.\par O \emph{dualismo} foi chamado tambem \emph{Manicheismo}, porque foi \emph{Manes} ou \emph{Manicheu} quem primeiramente o ensinou. --- Nos tempos modernos foi resuscitado pelo sceptico \emph{Bayle}, ensinando que o \emph{dualismo}, embora seja absurdo \emph{a-priori}, comtudo póde provar-se \emph{a-posteriori}.\par O \emph{dualismo} é falso pelas seguintes razões:\par \emph{a}) \emph{Baséa-se num falso fundamento}. Porquanto, admitte que existe um principio inteiramente mau; porque, aliás, não poderia explicar-se a existencia do mal. Ora esta razão é falsa; porque o mal, derivando das causas creadas, é uma consequencia do \emph{bem limitado}.\par \emph{b}) \emph{Um principio inteiramente mau não póde existir}. Porquanto, esse principio --- ou é privado de toda a perfeição e entidade, --- ou falta-lhe apenas a bondade \emph{moral}, --- Na primeira hypothese, um ente inteiramente mau seria \emph{nada} e o \emph{nada} não existe. --- Na segunda, um principio inteiramente mau seria, ao mesmo tempo, infinito e finito. Seria \emph{infinito}, porque o systema suppõe-no tal. Seria \emph{finito}, porque faltar-lhe-hia uma perfeição.\par \emph{c}) \emph{O dualismo perverte a moral}. Com effeito, se todo o \emph{bem} devesse attribuir-se a um principio \emph{bom}, e todo o \emph{mal} a um principio \emph{mau}, o homem não poderia dizer-se e não seria causa do bem nem do mal, e por isso não seria merecedor de premio nem de castigo.\par \emph{d}) \emph{Esta hypothese não explica a existencia do bem e do mal}. Porquanto --- ou estes dois principios seriam eguaes, --- ou o principio \emph{bom} seria mais forte que o \emph{mau}, --- ou este mais forte que aquelle. Na primeira hypothese, as forças dos dois principios paralysar-se-hiam, e assim não existiria, sobre a terra, nem o bem nem o mal. Na segunda, não haveria mal algum. Na terceira, não haveria bem algum. Ora estas consequencias são falsas.\par \emph{Bayle} diz que os dois principios fizeram um pacto, para se não transtornarem mutuamente nas suas obras.\par Mas esta explicação é absurda. Se assim fosse, esses dois principios seriam fracos, e nenhum d'elles seria infinitamente perfeito. --- De mais, esse pacto ou seria bom ou mau. Se fosse bom, não poderia ser acceite pelo principio \emph{mau}; pois este, sendo inteiramente \emph{mau}, não póde fazer nem tolerar bem algum. Se fosse mau, não poderia ser acceite pelo principio \emph{bom}; pois este, segundo o \emph{dualismo}, não póde fazer nem tolerar mal algum.\par b) O \emph{polytheismo} é o erro que admitte varios \emph{deuses}.\par Este erro, que derivou sobre tudo da corrupção da vontade, não foi constante nem universal. --- Não foi \emph{constante}; porque não existia antes do \emph{diluvio}, e porque a primeira crença dos povos foi o \emph{monotheismo} (Cf. Hooke, \emph{Princip. rélig. natural. et révél.,} --- Frank, \emph{Études Orientales}). --- Não foi \emph{universal}; porque, além de muitos philosophos (como \emph{Anaxagoras}, \emph{Socrates}, \emph{Platão}, \emph{Aristoteles}, \emph{Cicero}), alguns povos, como o \emph{d'Israel}, professavam o \emph{monotheismo}, e muitos outros, embora admittissem deuses, comtudo reconheciam \emph{Jupiter}, como pae e senhor de todos os deuses.}.

%%>>> art-41

\section*{ARTIGO III.}
\addcontentsline{toc}{section}{ARTIGO III.}

{\small\noindent\textsc{Summario}: --- Intelligencia de Deus, sua natureza e objecto. --- Vontade de Deus, sua natureza e objecto. --- Vida divina, sua existencia e natureza.\par}

\pg{462} \textbf{Attributos de Deus positivos.}

38. \textbf{Attributos de Deus positivos.} --- Attributos de \textsc{Deus} \emph{positivos} são os que \emph{affirmam} de \textsc{Deus} as perfeições dos seres creados, mas isentas de toda a imperfeição ou de todo o limite. --- Os principaes são: a \emph{verdade}, a \emph{bondade} e a \emph{belleza} \footnote{Depois dos attributos \emph{negativos} de \textsc{Deus}, devemos considerar os attributos \emph{positivos}. E primeiramente consideraremos os attributos \emph{transcendentaes}, que \emph{analogicamente} convêm a \textsc{Deus} e ás creaturas. Como dissemos na Ontologia, attributos \emph{transcendentaes} são os que convêm a \emph{todo o ente}, emquanto \emph{tal}, e que, embora não sejam realidades distinctas do proprio ente, comtudo são explicações mais claras do mesmo, porque exprimem o que o ente não exprime. --- São quatro: \emph{unidade}, \emph{verdade}, \emph{bondade}, \emph{bellezza}. --- Tendo tratado da \emph{unidade}, resta-nos tratar dos outros attributos.}.

39. \textbf{Verdade de Deus.} --- \emph{Verdade} de \textsc{Deus} é o attributo, que consiste na \emph{conformidade}, ou \emph{adequação}, entre o acto da intelligencia divina e o ser divino \footnote{S. Thomaz trata d' este assumpto --- na \emph{Sum}. \emph{Th}., p. I, q. 16, --- c. \emph{Gentes}, l. I, c. 60 e seg., --- \emph{De verit}., q. I, --- I \emph{Sent}. Dist. 19, q. 5, a. a. 2 e 3.}.

40. \textbf{Deus é verdadeiro.} --- Um ente é \emph{verdadeiro}, quando ha conformidade entre o seu ser e o acto da intelligencia, á qual de algum modo se refere. Ora entre o ser de \textsc{Deus} e o acto da sua intelligencia não só ha \emph{conformidade}, mas absoluta \emph{identidade}; porque, em \textsc{Deus}, a substancia e os actos das faculdades são uma e a mesma coisa. Logo \textsc{Deus} é \emph{verdadeiro} \footnote{A verdade, por ser uma adequação entre a intelligencia e o objecto, deve encontrar-se em \textsc{Deus}, cuja intelligencia é uma e a mesma coisa com a essencia. Na sua essencia, como veremos, \textsc{Deus} conhece todas as coisas. --- Nem se diga que a verdade, consistindo \emph{na composição} e na \emph{divisão}, não pode encontrar-se em \textsc{Deus}, cuja intelligencia é \emph{simplicissima}, como é a sua essencia. O principio --- a \emph{verdade consiste na composição e na divisão} --- verifica-se unicamente com relação á verdade, que é propria da intelligencia \emph{humana}, e não pode applicar-se á intelligencia \emph{divina}. Com effeito, quando se diz que a verdade consiste na adequação entre a intelligencia e o objecto, --- não devemos entender que o \emph{acto} da intelligencia deve por tal modo egualar o \emph{objecto}, que, se este fôr composto e material, tambem composto e material deva ser aquelle (o que é absurdo), --- mas devemos entender que o \emph{verbo} da intelligencia deve por tal forma corresponder ao seu \emph{objecto}, que tudo o que o \emph{verbo} exprime se deve encontrar realmente no \emph{objecto}, o qual pode ser simples ou composto, espiritual ou material. Portanto a \emph{qualidade}, isto é, a \emph{simplicidade} ou a \emph{composição} do \emph{acto intellectual} não entra na razão ou essencia da \emph{verdade}, mas refere-se exclusivamente ao principio intelligente; e por isso, se a intelligencia do homem, por ser\pgn{463}a menos perfeita, não pode alcançar a verdade se não por meio da \emph{composição} e da \emph{divisão} (juizo \emph{affirmativo} ou \emph{negativo}), a intelligencia de \textsc{Deus}, por ser infinitamente perfeita e simplicissima, como a sua essencia, com um acto unico e simplicissimo alcança toda a verdade, quer esta se refira á essencia divina, quer se refira ás coisas creadas, simples ou compostas, espirituaes ou materiaes. (\emph{Cf}. \emph{C}. \emph{Gentes}, l. I. c. 59).}.

\pg{463} 41. \textbf{Deus é a verdade.} --- Consistindo a verdade na conformidade entre o acto da intelligencia e o ser, \textsc{DEUS} será a propria \emph{verdade}, se fôr o seu entender e o seu proprio ser. Ora \textsc{Deus} é, effectivamente, o seu entender e o seu proprio ser. Logo \textsc{Deus} é a \emph{verdade}. --- \textsc{Deus} é a verdade \emph{summa} e \emph{infinita}, porque summa e infinitamente perspicaz é a sua intelligencia, e summa e infinitamente intelligivel é o seu ser \footnote{Em \textsc{Deus} não só ha a verdade, mas Elle é a propria verdade. Por quanto, tudo o que está em Deus não é \emph{partecipado} nem \emph{accidental}, mas é \emph{essencial} e \emph{substancial}, é o proprio ser divino, é o proprio Deus. Logo Deus, assim como é a sua essencia, o seu ser, assim tambem é a sua verdade. Diz S. Thomaz: « De Deo nihil \emph{participative} dici potest, cum sit suum esse, quod nihil participat. Sed veritas est in Deo. Si igitur non dicatur \emph{participative}, oportet quod dicatur \emph{essentialiter}. \textsc{DEUS} ergo est sua veritas ». (\emph{C}. \emph{Gentes}, l. I, c. 60). --- A verdade divina é tão pura, que não tolléra sombra de falsidade ou de engano. « A verdade, diz o mesmo Angelico, não é compativel com a falsidade, como a candura não é compativel com a côr preta. Ora Deus não só é verdadeiro, mas é a propria verdade. Logo em Deus não pode haver falsidade » (\emph{Ib}., c. 61).}.

42. \textbf{Deus é a verdade na ordem ontologica e na ordem logica.}

a) \textsc{Deus} É A VERDADE NA ORDEM ONTOLOGICA. --- \textsc{Deus} será a verdade na ordem \emph{ontologica}, se a sua essencia ou substancia se conformar com a sua intelligencia; porque a verdade \emph{ontologica} consiste na conformidade do objecto com a intelligencia. Ora a essencia de \textsc{Deus}, não só se conforma com a sua intelligencia, mas identifica-se com ella, e \textsc{Deus} é a sua essencia. Logo Deus é a verdade na ordem \emph{ontologica}.

b) \textsc{Deus} É A VERDADE NA ORDEM LOGICA. --- \textsc{Deus} será a verdade na ordem \emph{logica}, se a sua intelligencia se conformar com a sua essencia; porque a verdade \emph{logica} consiste na conformidade da intelligencia com o objecto. Ora, como dissemos, a intelligencia de \textsc{Deus} não só se conforma com a sua essencia, mas identifica-se com ella, e \textsc{Deus} é a sua intelligencia. Logo \textsc{Deus} é a verdade na ordem \emph{logica} \footnote{A verdade, como dissemos, pode considerar-se no \emph{objecto} e na \emph{intelligencia}. Considerada no \emph{objecto}, é a conformidade d'este com a intelligencia (\emph{veritas rei}); --- considerada na intelligencia, é a conformidade d'esta com o objecto (\emph{veritas intellectus}). Em Deus ha a verdade do \emph{objecto} ou \emph{ontologica}, e ha a verdade da \emph{intelligencia} ou \emph{logica}. Diz S. Thomaz: « Veritas\pgn{464}invenitur \emph{in intellectu}, secundum quod apprehendit \emph{rem ut est}; et \emph{in re}, secundum quod habet esse \emph{conformabile intellectui}. Hoc autem maxime invenitur in Deo. Nam esse suum non solum est conforme suo intellectui, sed etiam est ipsum suum intelligere..., et ipse est suum esse et intelligere. Unde sequitur quod non solum in ipso sit veritas, sed quod ipse sit ipsa summa veritas » (\emph{Sum}. \emph{Th}., p. I, q. 16, a. 5).}.

\pg{464} 43. \textbf{A verdade ontologica de Deus é a causa efficiente da verdade ontologica das creaturas.} --- A verdade \emph{ontologica} de \textsc{Deus} será a causa \emph{efficiente} da verdade \emph{ontologica} das creaturas, se o ser d'estas for effeito e expressão da Essencia divina. Ora o ser das creaturas, tendo sido produzido por \textsc{Deus} e em conformidade com as idéas archetypas, que são imitações diversas da Essencia infinita, é, na sua physica realidade, expressão e effeito da mesma Essencia. Logo a verdade \emph{ontologica} de \textsc{Deus} é a causa \emph{efficiente} da verdade \emph{ontologica} das creaturas \footnote{A Essencia divina é uma simplicissima unidade, mas é imitavel e participavel, por infinitos e diversos modos, pelas creaturas. « Cada creatura, diz S. Thomaz, possue a propria especie, emquanto de algum modo participa da semelhança da Essencia divina ». (\emph{Sum}. \emph{Th}., p. 1, q. 15, a. 2). Por isso cada realidade creada é uma \emph{expressão}, ou \emph{semelhança}, da unica Essencia divina. --- E' uma \emph{expressão}, e tambem um \emph{effeito}. Diz o mesmo santo Doutor: « Um ente, que possue o ser \emph{por essencia}, é necessariamente causa de todos os entes, que o possuem por \emph{participação}... Mas Deus possue o ser pela sua essencia, porque é o proprio ser; e todos os outros entes \emph{só} o possuem por participação, porque o ente, que é o proprio ser, é unico. Portanto, Deus é causa do ser de todos os outros entes » (\emph{C}. \emph{Gentes}, l. II, c. 15). --- Sendo uma expressão e um effeito da Essencia divina, o ser das creaturas não pode deixar de corresponder ás ideas archetypas da intelligencia divina. Por isso, o ser das creaturas refere-se \emph{necessariamente} á intelligencia divina, emquanto não pode deixar de corresponder ás ideas archetypas, e só \emph{accidentalmente} á intelligencia creada, emquanto \emph{pode} ser conhecido por ella; e a verdade \emph{ontologica} das creaturas é \emph{necessaria} com relação a Deus, e é \emph{contingente} com relação a nós. --- Como tambem essa verdade é \emph{unica} e \emph{multipla}. E' \emph{unica} com relação a Deus: porque unica é a Essencia divina, que pode ser e é diversamente imitada e participada. E' \emph{multipla} com relação ás creaturas; porque, consistindo a verdade \emph{ontologica} na conformidade do ente com a intelligencia, são tantas as verdades, quantas são as entidades.}.

44. \textbf{A verdade logica de Deus é a causa exemplar da verdade logica da intelligencia humana.} --- A verdade \emph{logica} da intelligencia humana depende das coisas, emquanto a nossa intelligencia, para alcançar e possuir essa verdade, deve conformar-se com as proprias coisas. Ora conformar-se com as coisas é conformar-se com as ideas archetypas da intelligencia divina, de que as coisas dependem essencialmente e que as coisas exprimem na sua realidade. Logo as ideas da nossa intelligencia emtanto \pg{465}  são dotadas de verdade \emph{logica}, emquanto imitam as idéas archetypas divinas, isto é, a verdade \emph{logica} de \textsc{Deus}. Logo a verdade \emph{logica} de \textsc{Deus} é a causa \emph{exemplar} da verdade \emph{logica} da intelligencia humana \footnote{« A verdade divina, diz S. Thomaz, é a medida de toda a verdade. Por quanto, a verdade da nossa intelligencia depende e é medida pela coisa, que está fóra da alma; pois dizemos que a nossa intelligencia é verdadeira, emquanto se conforma com a coisa. Mas a verdade da coisa é medida pela intelligencia divina, que é a causa de todos os seres creados, como a verdade das coisas artificiaes é medida pela arte do artista. Sendo, pois, \textsc{Deus} a primeira intelligencia e o primeiro intelligivel, é forçoso concluir que a sua verdade é a medida da verdade da intelligencia creada. Por isso a verdade divina é a \emph{primeira} verdade » (\emph{C}. \emph{Gentes}, l. 1, c. 62). --- Portanto a verdade \emph{logica}, propria da nossa intelligencia, pode dizer-se, ao mesmo tempo, \emph{unica} e \emph{multiplice}. E' \emph{unica} com relação a \textsc{Deus}, emquanto as nossas ideas são a expressão ou a semelhança, ainda que imperfeita, da unica e substancial verdade da intelligencia divina. E' \emph{multiplice} com relação ás \emph{formas} da intelligencia humana, que existem em muitas e distinctas intelligencias e representam muitos e diversos objectos.}.

\textbf{45. A verdade divina é a causa final de toda a verdade creada.} --- \textsc{Deus} não só é o principio, mas tambem é o fim de toda a creação; de modo que toda a perfeição creada não existe senão para imitar e revelar, embora de um modo inadequado, a correspondente perfeição divina. Ora a verdade creada é um effeito, e a sua existencia não pode ter e não tem outro fim, senão o de imitar e manifestar a verdade divina. Logo a verdade divina é a causa \emph{final} de toda a verdade creada \footnote{A verdade creada é, como dissemos, \emph{ontologica} e \emph{logica}, conforme se considera no \emph{objecto} ou na \emph{intelligencia}. Ora --- a verdade \emph{ontologica}, que é o proprio \emph{objecto} emquanto se conforma com a intelligencia, imita e manifesta a verdade \emph{ontologica} de \textsc{Deus}, porque todo o ente participado revela o ente por essencia, do qual deriva; --- e a verdade \emph{logica}, que é o verbo da intelligencia emquanto se conforma com o objecto, revela a verdade da intelligencia divina. Por isso a verdade creada, emquanto imita e revela a verdade divina, tende para esta, como para o seu fim natural; e d'este modo a verdade divina é a causa \emph{final} da verdade creada.}.

\textbf{46. Bondade de Deus.} --- \emph{Bondade} de \textsc{Deus} é o attributo, que consiste na \emph{conveniencia} da infinita perfeição de \textsc{Deus} com a sua santissima vontade. --- A \emph{bondade}, pois, é a propria entidade, a propria perfeição do ente, emquanto \emph{convem} á vontade, e attrahe-a ou pode attrahil-a \footnote{S. Thomaz trata da \emph{bondade} de \textsc{Deus} --- na \emph{Sum}. \emph{Th}. p. I, q. 6 --- c. \emph{Gentes}, l. I, c. 37 e seg. --- \emph{De verit.}, q. 21.}.

\textbf{47. Deus é bom.} --- Um ente é e diz-se \emph{bom}, emquanto é dotado de perfeição, que tem a força de attrahir a vontade.

\pg{466} Ora Deus contém em si toda a perfeição, pela qual é objecto de um amor sem limites da parte da sua vontade; pois todo o ente não pode deixar de amar a propria perfeição. Logo \textsc{Deus} é \emph{bom}. --- \textsc{Deus} é \emph{summa} e \emph{infinitamente bom}, porque \emph{summa} e \emph{infinita} é a sua entidade ou perfeição \footnote{Deus é \emph{bom}, e \emph{summa} e \emph{infinitamente} bom, porque é \emph{perfeito}, e \emph{summa} e \emph{infinitamente} perfeito. A bondade está na razão \emph{directa} da perfeição; porque a bondade não é senão a perfeição, emquanto se refere á vontade. Diz o P. Lessio; « Imagina uma infinita luz, um infinito poder, uma infinita sabedoria, uma infinita belleza, uma infinita doçura, uma infinita alegria, uma infinita riqueza, uma infinita gloria, uma infinita majestade, uma infinita santitade. Pois bem, a Bondade divina contém todas estas perfeições como um oceano infinito de todos os bens. E não só contem todo o bem, mas contém-no de um modo infinitamente elevado; porque não o contem por muitas e diversas formas ou qualidades, mas por uma simplicissima e elevadissima forma, que é a sua essencia. Este modo de conter as perfeições, torna estas infinitamente mais excellentes, do que se existissem distinctas uma de outra. (\emph{De perfectione divina}, l. 7, e. 1, n. 2).}.

48. \textbf{Deus} é a \textbf{bondade}. --- A \emph{bondade} de \textsc{Deus} é a sua perfeição; e por isso tal é a natureza da sua bondade, qual é a natureza da sua perfeição. Ora a perfeição, em \textsc{Deus}, não é uma coisa accidental, addicionada á essencia, mas é uma coisa substancial, é a divina essencia; porque toda a perfeição é um ente, e a essencia divina é o proprio ser. Logo \textsc{Deus}, assim como é a sua essencia, assim tambem é a sua \emph{bondade} \footnote{Se \textsc{Deus} não fosse a sua bondade, isto é, se não fosse bom pela sua \emph{essencia}, seria bom \emph{por participação}. Mas então, como o ente, que é bom \emph{por participação}, presuppõe outro ente, anterior a si, do qual recebeu a bondade, Deus não seria a Causa primeira, o Motor immovel, o Ente absoluto e infinito: o que é absurdo. Logo Deus é bom pela sua essencia, isto é, é a propria bondade. --- Se Deus é a propria Bondade, é claro que exclue absolutamente todo e qualquer \emph{mal}, que é opposto ao \emph{bem}. Alem d'isso, se o mal podesse existir em Deus, existiria como uma coisa \emph{essencial}, seria a propria \emph{essencia} de Deus, porque tudo o que está em Deus é o proprio Deus: o que é absurdo, porque o mal, sendo privação do bem, é nada, e o nada não pode constituir nenhuma essencia, muito menos a perfeitissima essencia divina. (Cf. c. \emph{Gentes}, l. I, c. 38 e 39).}.

49. A \textbf{bondade} divina é a causa efficiente de toda a bondade \textbf{creada}. --- A bondade divina será a causa \emph{efficiente} de toda a bondade creada, se as perfeições, de que são dotadas as creaturas, derivarem da unica e simplicissima bondade de \textsc{Deus}. Ora é certo que todas as perfeições, de que são dotadas as creaturas, derivam da unica e simplicissima bondade de \textsc{Deus}; porque as perfeições das creaturas são finitas, contingentes, mutaveis, e por isso derivam da perfeição divina, que é infinita \pg{467} nita, necessaria, immutavel. Logo a bondade divina é a causa \emph{efficiente} de toda a bondade creada \footnote{Toda a bondade das creaturas deriva, como da causa \emph{efficiente}, da bondade de Deus. Dizem os escholasticos: \emph{quod est maximum in unoquoque genere est causa aliorum quae sunt in illo genere}. --- Mas que differença entre a bondade de Deus e a das creaturas! A bondade de Deus é imparticipada e fonte inexhaurivel de todo o bem; a bondade das creaturas é participada e deriva de Deus. --- A bondade de Deus é illimitada, porque não ha quem a possa limitar; a bondade das creaturas é limitada, porque é recebida numa essencia finita. --- A bondade de Deus é substancial e identifica-se com a sua essencia; a bondade das creaturas é uma forma accidental e addicionada á essencia. --- A bondade de Deus é immutavel, não pode estar sujeita ao augmento nem á diminuição; a bondade das creaturas é capaz de augmento até ao infinito e de diminuição até ao nada. Por causa d'esta excellencia da bondade divina sobre a creada disse \textsc{JESUS CHRISTO}: « Só Deus é bom: \emph{Nemo bonus nisi solus Deus} » (Luc., XVIII, 19).}.

\textbf{50. A bondade divina é a causa exemplar da bondade creada.} --- O ente, que possue uma determinada qualidade por participação, não se diz dotado d'essa qualidade, se não emquanto manifesta uma certa \emph{semelhança} com o ente, que possue a mesma qualidade por essencia; assim como o ferro não se diz igneo, senão emquanto tem uma certa semelhança com o fogo. Ora \textsc{Deus} é bom por essencia, e todos os seres creados são bons por participação. Logo todo o ente creado diz-se \emph{bom}, emquanto tem uma certa semelhança com a bondade de Deus. Logo a bondade divina é a causa \emph{exemplar} da bondade creada \footnote{Cf. \emph{c}. \emph{Gentes}, l. 1, c. 40. --- A bondade divina é a causa \emph{exemplar}, mas não é a causa \emph{formal} da bondade creada, porque a bondade das creaturas não é uma e a mesma coisa com a bondade divina. A bondade creada é \emph{distincta} da bondade divina. --- Por isso a bondade creada é \emph{unica} e \emph{multiplice}. E' \emph{unica}, emquanto se refere a Deus, como ao seu primeiro principio \emph{efficiente} e \emph{exemplar}. E' \emph{multiplice}, emquanto a bondade ou perfeição creada se acha multiplicada tantas vezes, quantas são as creaturas.\par Advertimos que as creaturas não participam da \emph{bondade} de \textsc{Deus} pelo mesmo modo, por que participam da sua \emph{verdade}. Com effeito, as creaturas participam da \emph{bondade} de \textsc{Deus} --- não só \emph{extrinsecamente}, emquanto derivam de \textsc{Deus}, --- mas tambem \emph{intrinsecamente}, emquanto possuem em si mesmas a forma de bondade, que é uma semelhança, embora infinitamente remota e deficiente, da propria bondade de \textsc{Deus}; ao passo que ellas se dizem \emph{verdadeiras}, --- não \emph{intrinsecamente}, pois a verdade reside, como no seu proprio sujeito, na \emph{intelligencia}, e não no \emph{objecto}, --- mas apenas \emph{extrinsecamente}, emquanto os seres creados, por isso mesmo que são creados, se referem, e de um modo \emph{necessario}, á intelligencia divina. (Cf. Sum. \emph{Th}., p. I, q. 6, a. 4).}.

\textbf{51. A bondade divina é a causa final da bondade creada.} --- A bondade divina será a causa \emph{final} da bondade creada, se \textsc{Deus}, infundindo a bondade nas creaturas, tiver tido o fim de

\pg{468} communicar e manifestar a sua bondade infinita. Ora tal foi, effectivamente, o fim do Creador. Logo a bondade divina é a causa \emph{final} da bondade creada \footnote{Não pode duvidar-se de que Deus creou o mundo com o fim de communicar a outros seres a propria bondade. Porquanto, todo o agente, quando opéra, tem em vista a propria bondade ou perfeição --- para a adquirir, se não a possue, como o agente creado, --- e para a communicar, se a possue, como o agente increado.}.

52. \textbf{Belleza} de \textbf{Deus}. --- \emph{Belleza} de \textsc{Deus} é o attributo, que consiste na infinidade das perfeições divinas, reduzidas a uma unica e simplicissima unidade, e sobremaneira agradaveis a quem as contempla \footnote{Como dissemos na \emph{Ontologia}, a \emph{belleza é a conveniente unidade na variedade}, \emph{que agrada a quem a vê: pulchritudo est multitudinis debitae debita unitas}, \emph{quae visa placet}. Na belleza devemos distinguir dois elementos --- o elemento \emph{objectivo} e o \emph{subjectivo}. O elemento \emph{objectivo}, que constitue propriamente a essencia da belleza, é a \emph{conveniente unidade na variedade}, isto é, a devida multiplicidade reduzida á devida unidade; o elemento \emph{subjectivo}, que não é propriamente \emph{constitutivo} mas \emph{effeito} da belleza, é o \emph{agrado}, que a propria belleza produz em quem a contempla. --- O elemento \emph{objectivo é composto}. \emph{A multiplicidade} ou \emph{variedade} é a \emph{materia} da belleza; a \emph{unidade} é a \emph{forma}, como diz S. Agostinho: « Omnis pulchritudinis \emph{forma} est \emph{unitas} » (\emph{Ep}. 18, al. 63, ad Coelest. n. 2). --- O elemento \emph{subjectivo} indica que a belleza, como a verdade e a bondade, é uma coisa \emph{relativa}. A verdade refere-se á faculdade cognoscitiva, a bondade á appetitiva, a belleza refere-se á cognoscitiva e á appetitiva; porque todos dizem que \emph{pulchrum est quod visum placet}, ora uma coisa \emph{vé-se} pela faculdade cognoscitiva, e \emph{agrada} á appetitiva. --- A belleza é um attributo \emph{transcendental}; porque é proprio de todo o ente, emquanto \emph{tal}. Todo o ente é \emph{multiplo} (real ou logicamente) e \emph{uno}, e a sua belleza é maior ou menor, conforme o maior ou menor numero de partes e conforme a maior ou menor harmonia e união, que existe entre ellas.}.

53. \textbf{Deus} é \textbf{bello}. --- Consistindo a belleza na multiplicidade das perfeições devidas, reduzida á conveniente unidade, \textsc{Deus} será \emph{bello}, se possuir todas as perfeições, que são proprias do Ente increado, e se todas essas perfeições forem reduzidas á devida unidade; de modo que essa unidade em tanta variedade possa agradar a quem a contempla. Ora \textsc{Deus} possue perfeições infinitas, e todas essas perfeições constituem a unica e simplicissima Essencia divina, e tal unidade numa variedade infinita não pode deixar de agradar immensamente a quem a contempla. Logo \textsc{Deus} é \emph{bello} \footnote{A conclusão é evidente, e prova-se tambem com argumentos \emph{a-posteriori}. O mundo, creado por Deus, é dotado de belleza, e é por isso que foi chamado \emph{cosmos} pelos gregos, e \emph{mundo} pelos latinos. Ora toda a perfeição do effeito deve preexistir na sua Causa efficiente. --- Além d'isso,\pgn{469}a belleza, que se admira nas creaturas, é mutavel e imperfeita. Ora tudo o que é mutavel e imperfeito suppõe e exige uma causa absolutamente immutavel e perfeita. Logo deve existir um Ente immutavel e perfeitamente bello, de que deriva toda a belleza, espalhada nas creaturas.}.

\pg{469} 54. \textbf{Deus é a sua belleza.} --- Todas as perfeições divinas são uma e a mesma coisa com a essencia divina; porque \textsc{Deus} é perfeito pela sua essencia. Ora a belleza é uma perfeição. Logo a belleza é uma e a mesma coisa com a essencia divina; e por isso \textsc{Deus}, assim como é a sua essencia, assim tambem é a sua \emph{belleza}. --- Além d'isso, consistindo a belleza divina numa infinita variedade de perfeições, reduzidas a uma simplicissima unidade, é evidente que \textsc{Deus} será a sua belleza, se fôr a sua perfeição e a sua unidade. Ora \textsc{Deus} é a sua perfeição e a sua unidade; porque tudo o que está em \textsc{Deus}, é \textsc{Deus}. Logo \textsc{Deus} é a sua \emph{belleza} \footnote{\textsc{Deus} é a propria belleza! Quem poderá comprehender e exprimir todo o encanto, toda a suavidade, que esta expressão encerra?... A belleza constitue a propria Essencia de \textsc{Deus}!... Os escriptores dizem que a belleza de Deus é \emph{summa}, \emph{infinita}, \emph{incorruptivel}, \emph{immutavel}, \emph{amabilissima}; mas essas palavras levam-nos ao perfeito conhecimento da belleza divina? Não. Só \textsc{Deus} contempla e comprehende a sua belleza infinita e eterna; só \textsc{Deus} ama a sua belleza com amor infinito e eterno; só \textsc{Deus} é infinitamente bemaventurado na contemplação e no amor da sua belleza... A creatura racional vê um pallido reflexo d'esta belleza infinita nas coisas visiveis, e d'esse reflexo, depurado das imperfeições, a intelligencia forma um ideal de belleza, que é superior a toda a belleza creada e representa, ainda que de um modo muito inadequado e deficiente, a ineffavel belleza divina. Esse ideal é sufficiente para alma se destacar de todos os encantos creados e se unir intimamente, pela fé e pelo amor, ao Summo Bem... Isto é muito, mas não é tudo. Deus quiz, na sua infinita misericordia, elevar o homem a uma ordem sobrenatural, que tem por fim a visão \emph{clara} e \emph{intuitiva} da propria essencia de Deus. \emph{Videbimus Deum sicuti est}! A nossa intelligencia verá, pela graça de Deus, a propria essencia divina como é em si mesma, e por isso infinitamente bella, a propria belleza! Não será uma visão \emph{comprehensiva} (porque o finito não pode comprehender o infinito), mas será verdadeira \emph{visão}. Esta visão transformará a alma e a tornará semelhante a Deus, e será principio de um amor e de uma felicidade sem limites! Felizes os que, conscientes da sua excelsa dignidade e do seu eterno destino, desprezam as bellezas caducas e miseraveis das creaturas, e collocam as suas esperanças e os seus amores na unica e immutavel Belleza divina!}.

\textbf{55. A belleza divina é a causa efficiente de toda a belleza creada.} --- A belleza divina será a causa \emph{efficiente} de toda a belleza creada, se esta derivar de \textsc{Deus}, emquanto e porque Elle é dotado de belleza infinita. Ora é certo que a belleza creada deriva de \textsc{Deus}, emquanto e porque Elle é dotado de belleza infinita; porque ninguem pode communicar a outros uma \pg{470} perfeição, senão emquanto e porque a possue. Logo a belleza divina é a causa \emph{efficiente} de toda a belleza creada \footnote{De \textsc{DEus}, como da causa \emph{efficiente}, deriva toda a belleza, como toda e qualquer qualidade das creaturas. --- A belleza creada é multipla e diversa, porque muitos e diversos são os entes creados. São bellos os seres \emph{inorganicos}, sobretudo os crystaes e as pedras preciosas; principalmente a \emph{luz}, que é esplendida em si e torna esplendidos os outros corpos, e por isso é imagem e symbolo da Belleza divina. --- Mais bellos são os seres \emph{vivos}; sobretudo o homem, em cuja alma espiritual \textsc{Deus} gravou a propria imagem e semelhança. --- Immensamente, ineffavelmente bella é a SS. Humanidade de N. S. \textsc{Jesus Christo}, pela integridade, regularidade e harmonia das suas feições immaculadas, por aquelle esplendor de magestade e de graça, com que a Divindade, que nella habita substancialmente, a circumda e glorifica. --- Mas de todas essas bellezas a fonte unica é a Belleza divina.}.

56. \textbf{A belleza divina é a causa exemplar de toda a belleza creada.} --- Todo o agente tende a imprimir e imprime, de algum modo, no effeito a semelhança da propria forma, por que opera. Ora Deus opera pela sua essencia, que é bella e é a propria belleza. Logo Deus, produzindo as creaturas, imprime nellas a semelhança da sua belleza. Logo a belleza divina é a causa \emph{exemplar} de toda a belleza creada \footnote{Deve haver uma \emph{analogia}, uma \emph{semelhança} entre a belleza divina e a belleza creada, como ha sempre analogia ou semelhança entre a causa e o effeito. --- Mas a \emph{analogia}, a \emph{semelhança} não é \emph{egualdade}. \textsc{Deus} não é causa \emph{univoca} das creaturas; por isso a forma das creaturas não pode ser \emph{especificamente egual} á forma do Creador. Portanto a belleza divina é infinitamente mais pura, mais elevada, mais attrahente. Ora, se a belleza das creaturas, embora imperfeita e caduca, é digna de estima e tanto prende o coração, de quanta estima não é digna e quanto não deve ser amada a infinita e purissima belleza de Deus?!}.

57. \textbf{A belleza divina é a causa final da belleza creada.} --- A belleza divina será a causa \emph{final} da belleza creada, se \textsc{Deus} dotou de belleza as creaturas, para que estas representassem a sua belleza infinita. Ora a intenção de \textsc{Deus}, na producção do mundo, foi effectivamente a de revelar nas visiveis perfeições das creaturas as proprias perfeições invisiveis. Logo a belleza divina é a causa \emph{final} da belleza creada \footnote{Se a belleza divina é a causa \emph{final} da belleza creada, esta belleza creada não pode nem deve considerar-se como o \emph{termo} e o \emph{fim ultimo} das nossas tendencias e dos nossos desejos, mas unicamente como o \emph{meio} e o \emph{degrau} para nos elevarmos á contemplação e ao amor da belleza divina. O espirito, que assim considera a belleza e, em geral, as perfeições das creaturas, não pode deixar de se unir intimamente ao infinito Bem.\par Concluimos este artigo e confirmamos as nossas considerações com as palavras do nosso Angelico Doutor: « Consideratio creaturarum animos hominum in amorem divinae bonitatis accendit. Quidquid enim bonitatis et perfectionis\pgn{471}in diversis creaturis particulariter distributum est, totum in Deo universaliter est adunatum, sicut in fonte totius bonitatis. Si igitur creaturarum bonitas, pulchritudo et suavitas sic animos hominum allicit; ipsius Dei fons bonitatis, rivulis bonitatis in singulis creaturis repertis diligenter comparatus, animos hominum inflammatos totaliter ad se trahet » (C. \emph{Gentes}, l. II, c. 2).}.

\pg{471}
